\documentclass[12pt,a4paper,openright,twoside]{report}

\usepackage[utf8]{inputenc}
\usepackage[a4paper,inner=3.5cm,outer=2.5cm]{geometry}

\usepackage[titletoc,title,toc,page]{appendix}
\usepackage{verbatim}
\usepackage{placeins}
\usepackage{hyperref}
\hypersetup{pdfborder=0 0 0}
\usepackage[italian]{babel}
\usepackage{tikz}
\usetikzlibrary{positioning, shapes.geometric, fit, backgrounds, arrows.meta}
\usepackage{parskip}

\usepackage{graphicx}
\usepackage{svg}
\usepackage{pdfpages}
\usepackage{blindtext}
\usepackage{chngcntr}
\counterwithin{table}{chapter}

\usepackage{amssymb}
\usepackage{amsmath}
\usepackage{fancyhdr}
\usepackage{indentfirst}
\usepackage{float}
\usepackage[capitalize,noabbrev]{cleveref}
\usepackage{soul}
\usepackage[font=footnotesize,labelfont=bf]{caption}

\usepackage{multirow}
\usepackage{hyphenat}
\hyphenation{mate-mati-ca recu-perare}

\usepackage{lscape} 

\usepackage{natbib}
\bibliographystyle{plain}
\setcitestyle{super,open={[},close={]}}

\newcommand{\rom}[1]{\uppercase\expandafter{\romannumeral #1\relax}}

\usepackage{listings}
\input{listings_highlights/sharpc-highlighting.tex}
\input{listings_highlights/rust-highlighting.tex}
\lstdefinestyle{listingsStyle}{
  numbers=none, basicstyle=\small, backgroundcolor=\color{white}
}

\addtolength{\skip\footins}{2pc plus 5pt}

\setcounter{tocdepth}{1}

\begin{document}

\includepdf[pages=-]{frontespizio/frontespizio.pdf}

\large

\pagenumbering{gobble}
\chapter*{Abstract}

Nelle reti blockchain decentralizzate la velocità complessiva del sistema è dettata dal suo nodo più lento. Per preservare un ambiente trustless, il protocollo impone a ciascun validatore la ri-esecuzione nativa di ogni transazione e il mantenimento di un database di stato dell'ordine dei terabyte. Questo approccio comporta una complessità asintotica $O(N)$ rispetto al numero di transazioni da processare e innesca il fenomeno dello State Bloat, allontanando progressivamente i dispositivi hardware consumer dalla partecipazione attiva alla rete. Questa tesi presenta ZeroNethermind, un Proof-of-Concept per trasformare il client Execution Layer Nethermind in un validatore Stateless per abbattere la barriera computazionale $O(N)$ attraverso l'adozione di un paradigma di verifica crittografica con complessità $O(1)$. Sfruttando l'efficienza delle Zero-Knowledge Proofs, ZeroNethermind disaccoppia il throughput transazionale dalle capacità hardware del validatore. I nodi cessano di ri-eseguire le transazioni, limitandosi a verificare matematicamente una prova crittografica compatta che certifica l'integrità del blocco. Delegando il massiccio carico di esecuzione a prover specializzati ad alte prestazioni e consentendo la verifica in frazioni di secondo su dispositivi consumer, ZeroNethermind stabilisce le fondamenta per un livello di esecuzione completamente stateless, garantendo una scalabilità massiva pur preservando la natura trustless e decentralizzata del protocollo.

\topmargin=-1cm
\tableofcontents
\thispagestyle{empty}
% \listoftables
% \thispagestyle{empty}
% \listoffigures
% \thispagestyle{empty}
\newpage

\pagenumbering{arabic}
\setcounter{chapter}{0}
\raggedbottom
\setcounter{page}{1}
\pagestyle{plain}


% ============================================================
%  CAPITOLO 1: INTRODUZIONE
% ============================================================
\chapter{Introduzione}\label{chap:introduzione}

La promessa fondamentale delle reti blockchain decentralizzate risiede nella democratizzazione della fiducia e nella resistenza sistemica alla censura. Questa \textit{libertà} algoritmica, tuttavia, si scontra con rigorose limitazioni fisiche e architetturali. Nell'ecosistema Ethereum, il mantenimento dell'integrità del network in un ambiente trustless richiede attualmente che ogni nodo validatore ri-esegua nativamente ogni singola transazione. Di conseguenza, il throughput complessivo della rete è vincolato dalle capacità computazionali del nodo più lento all'interno del sistema.

Questo collo di bottiglia architetturale impone un costo di ri-esecuzione pari a $O(N)$, proporzionale al volume delle transazioni. Aumentare il Gas Limit per processare un volume maggiore di transazioni comporterebbe l'inevitabile esclusione dell'hardware di livello consumer. A questa problematica si sovrappone la crisi architetturale nota come State Bloat, l' espansione dello stato globale. È fondamentale precisare che il principale collo di bottiglia nell'elaborazione e validazione dei blocchi su Ethereum non è di natura computazionale (CPU-bound), bensì legato all'accesso ai dati (I/O-bound). L'esecuzione di una singola transazione all'interno dell'EVM richiede continue e frammentate operazioni di lettura e scrittura sul database di stato, il quale deve tracciare accuratamente la totalità dei saldi degli account e degli slot di memoria degli smart contract.

Attualmente, questo strato dati impone ai nodi il mantenimento locale di archivi nell'ordine dei terabyte. Con la crescita ininterrotta dello stato, i requisiti hardware si inaspriscono: il limite non è più dettato dalla capacità di archiviazione, ma dalla necessità di disporre di memorie con prestazioni elevatissime in termini di IOPS (Input/Output Operations Per Second) per sostenere il ritmo di esecuzione. Tale deriva spingerebbe la validazione verso una rete dominata da data center centralizzati, minacciando il fondamento stesso della decentralizzazione e della resistenza alla censura.

Con la transizione dalla fase architetturale di "The Merge" alla roadmap futura di "The Verge"\cite{verge_2024}, l'ecosistema Ethereum sta attivamente puntando a mitigare questa problematica. Il suo obiettivo a lungo termine è raggiungere la \textit{statelessness}, abbassando i requisiti hardware al punto da permettere la validazione completa della chain su dispositivi a bassissimo consumo, come smartphone o smartwatch.

Storicamente, l'evoluzione della crittografia ha tracciato la roadmap per la sicurezza delle infrastrutture digitali mondiali: l'introduzione delle funzioni di hash, della crittografia a chiave asimmetrica e delle firme digitali ha posto le fondamenta per l'autenticazione e il trasferimento di valore senza intermediari. L'adozione della crittografia Zero-Knowledge (ZK) rappresenta il passo evolutivo successivo.

Le architetture blockchain tradizionali sono intrinsecamente "simmetriche": ogni nodo è costretto a duplicare lo stesso identico sforzo computazionale del nodo che ha prodotto il blocco. L'applicazione delle ZK-Proofs (ZKP) impone un cambio di paradigma radicale: lo sforzo per generare il blocco e la relativa prova crittografica resta oneroso, ma la successiva verifica diviene computazionalmente irrisoria. In un ecosistema basato su ZKP, la validazione computazionale di un blocco contenente diecimila transazioni richiede il medesimo tempo necessario per verificare un blocco contenente una singola transazione.

Questo principio si concretizza nel modello topologico "Beast vs. Fort"\cite{lean_ethereum}. Da un lato, le infrastrutture ``Beast Mode`` (Provers/Builders) centralizzate e dotate di cluster hardware ad altissime prestazioni generano i blocchi e le relative prove crittografiche. Dall'altro lato, i validatori in ``Fort Mode`` operano su hardware di fascia consumer, rimuovendo il peso dell'esecuzione e limitandosi a certificare la validità matematica della prova e degli header dei blocchi, garantendo una sicurezza incompromissibile a difesa della rete.

L'obiettivo primario di questo lavoro di tesi è l'ingegnerizzazione di \textbf{ZeroNethermind}: un Proof-of-Concept ideato per trasformare il client Execution Layer Nethermind in un nodo ``Fort Mode``, finalizzato a validare i blocchi della Mainnet (L1) utilizzando ZKP, svincolandosi totalmente dalla necessità di mantenere un database di stato locale.

Il presente elaborato apporta i seguenti contributi al panorama dei client Ethereum:
\begin{itemize}
    \item \textbf{Implementazione del Paradigma "Double Zero":} Dimostrazione pratica della fattibilità di un \textit{validatore stateless} (zero-state) che integra zero-knowledge per sostituire la ri-esecuzione all'interno di un client production-ready.
    \item \textbf{Interoperabilità Cross-Language:} Sviluppo di un bridge FFI (Foreign Function Interface) ad alte prestazioni tra l'ambiente C\# (.NET) e le librerie crittografiche in Rust.
    \item \textbf{Validazione Live L1:} Intercettazione del Data Flow dell'Engine API per validare in tempo reale i blocchi in cima alla Mainnet demandando la generazione delle prove a infrastrutture esterne.
    \item \textbf{Analisi Quantitativa e Prestazionale:} Valutazione empirica dei tempi di verifica crittografica rispetto all'esecuzione nativa EVM e analisi dell'abbattimento del resource footprint in una Ethereum devnet privata.
\end{itemize}


Il perimetro della ricerca si focalizza sull'Execution Layer (EL), sulle dinamiche di intercettazione dell'Engine API e sulla fattibilità architetturale dell'approccio stateless. Non rientra negli scopi di questo lavoro l'ingegnerizzazione dell'infrastruttura di proving (Beast Mode), né la generazione locale delle prove crittografiche. Inoltre, la formulazione teorica estesa dei circuiti aritmetici sottostanti i protocolli ZK viene trattata solo in funzione della loro applicazione pratica per la verifica.

Il documento è strutturato come segue. Il \texttt{capitolo~\ref{chap:stato_arte}} presenta lo Stato dell'Arte, analizzando l'architettura di Ethereum, la dicotomia di scalabilità tra Layer~1 e Layer~2, il paradigma dei client stateless e i fondamenti della crittografia Zero-Knowledge, concludendo con una rassegna dei lavori correlati e delle proposte di aggiornamento del protocollo. Il \texttt{capitolo~\ref{chap:zeronethermind}} descrive la progettazione di ZeroNethermind: il paradigma ``Double Zero'', l'architettura a tre strati e il flusso di dati modificato rispetto alla validazione standard. Il \texttt{capitolo~\ref{chap:implementazione}} documenta l'implementazione condotta attraverso tre distinte fasi incrementali. Il \texttt{capitolo~\ref{chap:validazione}} presenta la validazione sperimentale su due fronti: una Demo live sulla Mainnet Ethereum e un ambiente controllato basato su devnet Kurtosis per l'analisi comparativa. Il \texttt{capitolo~\ref{chap:conclusioni}} chiude l'elaborato con le conclusioni, le limitazioni del PoC e le prospettive di sviluppo futuro nel contesto della roadmap di Ethereum.

% ============================================================
%  CAPITOLO 2: STATO DELL'ARTE
% ============================================================
\chapter{Stato dell'Arte}\label{chap:stato_arte}

Il presente capitolo ha l'obiettivo di delineare il contesto tecnologico e scientifico in cui si inserisce questo lavoro di tesi, fornendo le basi necessarie per comprendere le sfide attuali e le soluzioni emergenti nell'ecosistema blockchain. Il capitolo è strutturato in quattro sezioni principali: la prima analizza l'Architettura di Ethereum, definendone i principi di base, i limiti prestazionali intrinseci e il funzionamento dei nodi; la seconda esplora la dicotomia della scalabilità (L1 vs L2); la terza analizza l'introduzione di Client Stateless all'interno del protocollo di Ethereum; la quarta sezione è dedicata alla crittografia Zero-Knowledge, il motore del cambio di paradigma da computazione a verifica; infine, l'ultima sezione passa in rassegna i lavori correlati e il panorama della ricerca corrente.

\section{Architettura di Ethereum}

Ethereum è la piattaforma blockchain \textit{general-purpose} più adottata a livello globale per l'esecuzione di \textit{smart contract}. Più che una semplice rete per il trasferimento di valore, Ethereum è stato concepito come un'infrastruttura su cui chiunque può costruire applicazioni decentralizzate (dApp) senza richiedere permessi (\textit{permissionless}). I tre pilastri fondamentali su cui si basa questa architettura sono:
\begin{itemize}
    \item \textbf{Trasparenza:} Tutto il codice e lo storico delle transazioni sono pubblici e verificabili crittograficamente in modo indipendente.
    \item \textbf{Resistenza alla censura:} Nessuna entità centrale, governo o azienda ha l'autorità o la capacità tecnica di bloccare una transazione valida o impedire l'accesso alla rete.
    \item \textbf{Controllo da parte degli utenti:} Tramite l'uso della crittografia a chiave asimmetrica, gli utenti mantengono la piena e unica sovranità sui propri fondi e sui propri dati (\textit{self-custody}).
\end{itemize}
La novità dirompente introdotta da Ethereum è stata la creazione del primo "computer globale decentralizzato". A differenza dei paradigmi informatici tradizionali centralizzati (es. servizi cloud gestiti da entità corporative), l'infrastruttura di Ethereum non è localizzata in un singolo luogo fisico. È una macchina virtuale distribuita inestensibile che non può essere spenta, resettata o distrutta. L'accesso richiede unicamente una connessione a Internet, e l'infrastruttura crittografica garantisce uno spazio di indirizzamento enorme: le chiavi private derivano da uno spazio di probabilità pari a $2^{256}$, da cui vengono estratti indirizzi a 160 bit, garantendo virtualmente spazio inesauribile per l'intera popolazione globale senza alcun rischio pratico di collisioni.

\subsection{Blockchain Trilemma}
Sebbene Ethereum funga da computer globale, esso è intrinsecamente più lento e costoso rispetto all'informatica centralizzata. Questa non è un'inefficienza ingegneristica casuale, ma una precisa scelta architetturale legata al "Blockchain Trilemma". Tale postulato afferma che una rete decentralizzata deve bilanciare tre proprietà — Scalabilità, Sicurezza e Decentralizzazione — potendone storicamente massimizzare solo due simultaneamente. 

\begin{figure}[H]
    \centering
    \includegraphics[width=1\textwidth]{img/trilemma.png}
    \caption{Blockchain Trilemma\cite{Network_2025}}
    \label{fig:trilemma}
\end{figure}

Ethereum ha sistematicamente privilegiato Sicurezza e Decentralizzazione.
I colli di bottiglia primari derivano da due fattori critici:
\begin{enumerate}
    \item \textbf{Ri-esecuzione $O(N)$:} Per garantire che il sistema sia \textit{trustless}, ogni singolo nodo validatore della rete deve ri-eseguire in modo indipendente ogni singola transazione per verificarne la validità. Questo lavoro massicciamente ridondante impone un costo asintotico $O(N)$ sull'esecuzione, rendendo l'intero network veloce unicamente quanto il suo "nodo più lento".
    \item \textbf{Latenza di Rete e Consenso:} Ethereum è una rete P2P distribuita globalmente. Per permettere a migliaia di nodi di ricevere i dati, validare i blocchi e raggiungere un accordo probabilistico o deterministico sullo stato della rete, il protocollo impone la rigida divisione del tempo in \textit{Slot} fissi di 12 secondi.
\end{enumerate}

\subsection{Lo Stato in Ethereum e la Funzione di Transizione}

Ethereum opera fondamentalmente come una complessa macchina a stati finiti deterministica basata su transazioni. Il concetto di "Stato" (\textit{State}) è vitale: esso è il database globale che contiene il quadro aggiornato della rete.

La letteratura distingue tra \textit{World State} (Stato Globale) e \textit{Account State} (Stato dell'Account). Il \textit{World State} mappa ogni indirizzo a 160 bit al suo rispettivo \textit{Account State}. Un \textit{Account State} è composto da quattro elementi: il \textit{nonce} (contatore delle transazioni), il \textit{balance} (saldo in ETH), il \textit{codeHash} (l'hash del bytecode del contratto, se presente) e lo \textit{storageRoot} (la radice dell'albero che mappa le variabili di stato interne al contratto).

La transizione da uno stato valido al successivo è governata dalla \textit{State Transition Function} (STF) della Ethereum Virtual Machine (EVM). Come definito formalmente nello Yellow Paper\cite{yellow_paper}, l'applicazione di una transazione $T$ allo stato pre-esistente $\sigma_t$ genera un nuovo stato valido $\sigma_{t+1}$:\begin{equation}\sigma_{t+1} \equiv \Upsilon(\sigma_t, T)\end{equation}

Per garantire l'integrità crittografica di questo immenso database (che oggi pesa svariate centinaia di gigabyte, superando il terabyte per i nodi \textit{Archive}), Ethereum non utilizza un database relazionale standard, bensì un albero crittografico chiamato \textbf{Merkle Patricia Trie} (MPT). L'MPT funge da database autenticato composto da quattro tipologie di nodi: \textit{Branch} (nodi di diramazione a 16 vie), \textit{Extension} (nodi che comprimono percorsi comuni), \textit{Leaf} (nodi foglia contenenti il valore finale) e \textit{Empty} (nodi nulli). I dati grezzi serializzati di questo albero vengono fisicamente persistiti su disco tramite motori di database chiave-valore ad alte prestazioni come LevelDB o RocksDB. Tuttavia, il costo computazionale e di I/O (Input/Output) di questa architettura è estremamente gravoso: per ogni operazione di persistenza e aggiornamento, l'EVM deve ricalcolare crittograficamente gli hash partendo dal nodo foglia modificato e risalendo l'intero albero fino alla radice. Man mano che l'MPT cresce in dimensione e profondità, questo meccanismo di autenticazione genera una severa amplificazione delle operazioni di lettura e scrittura su disco (\textit{read/write amplification}), arrivando a pesare fino a 64 accessi al database per una singola operazione nel caso peggiore. Questo rende la ri-esecuzione $O(N)$ delle transazioni un problema prevalentemente \textit{I/O-bound} (limitato dalla latenza e dall'accesso sequenziale/casuale allo storage) e non puramente \textit{CPU-bound}. L'hash della radice del nuovo albero globale (\textit{State Root}) viene infine calcolato e inserito nell'intestazione di ogni nuovo blocco per sigillarne la validità.

\subsection{Gas Limit and Bottlenecks}
Il \textit{Gas} è l'unità di misura dello sforzo computazionale richiesto per eseguire un'operazione nell'EVM. Poiché Ethereum impone la re-esecuzione $O(N)$ su tutti i nodi, la rete deve autoproteggersi dal problema dell'arresto (Halting Problem) e dal consumo illimitato di risorse. Questo avviene tramite il \textit{Block Gas Limit}, un tetto massimo assoluto alla quantità di calcoli effettuabili in un singolo blocco (attualmente intorno ai 36 milioni di gas, in lenta evoluzione e oggetto di intensi dibattiti per aumenti fino a 60 milioni).

Aumentare "semplicemente" il Gas Limit comporterebbe un incremento drastico del carico computazionale, della larghezza di banda e dei requisiti di RAM e storage (SSD NVMe). L'aggravarsi dello \textit{State Bloat} — l'accumulo permanente di dati di stato — e l'aumento dell'onere di validazione spingerebbero inevitabilmente l'hardware di livello consumer fuori dalla rete, delegando il consenso a data center professionali e distruggendo la neutralità credibile dell'infrastruttura.

\subsection{Execution Layer (EL) vs Consensus Layer (CL)}

A seguito dell'aggiornamento "The Merge" (transizione al meccanismo Proof-of-Stake), l'architettura dei client Ethereum ha subìto una profonda refattorizzazione, separando le responsabilità in due moduli software eseguiti in tandem:
\begin{itemize}
    \item \textbf{Execution Layer (EL):} Questo strato (es. Nethermind, Geth) è il "motore" computazionale. L'EL mantiene l'MPT del \textit{World State}, gestisce la \textit{mempool} locale delle transazioni non confermate ed esegue nativamente l'EVM per applicare la funzione di transizione di stato.
    \item \textbf{Consensus Layer (CL):} Noto anche come \textit{Beacon Node} (es. Lighthouse, Prysm), questo strato è agnostico rispetto all'esecuzione degli \textit{smart contract}. Il suo compito è governare il protocollo di consenso Proof-of-Stake: coordina i turni dei validatori, propaga i blocchi tramite rete gossip, aggrega le firme (\textit{attestations}) e applica la regola di scelta del fork (\textit{Fork Choice Rule}) per stabilire la catena canonica
\end{itemize}

Questi due livelli comunicano in modo bidirezionale, sicuro e standardizzato tramite una suite di interfacce JSON-RPC nota come \textbf{Engine API}.

\subsection{Flusso di Validazione di un Blocco e l'Engine API}\label{sec:validazione_standard}

Il processo di validazione di un nuovo blocco rappresenta il momento in cui l'onere computazionale ricade fisicamente sul nodo. Le dinamiche di comunicazione tra EL e CL sono il fulcro su cui opera ZeroNethermind. Il ciclo di vita segue questi step architetturali:

\begin{figure}[H]
    \centering
    \includegraphics[width=\textwidth]{img/basic_validation_flow.png}
    \caption{Flusso di Validazione Standard}
\end{figure}

\begin{enumerate}
    \item \textbf{Mempool e Costruzione:} Le transazioni inviate dagli utenti popolano la \textit{mempool} dell'Execution Layer. Il nodo selezionato come \textit{Proposer} per lo slot corrente aggrega le transazioni in un \textit{ExecutionPayload} e lo incapsula in un \textit{Beacon Block}.
    \item \textbf{Propagazione (CL):} Il blocco Beacon viene diffuso attraverso la rete P2P del Consensus Layer.
    \item \textbf{Ricezione e Interrogazione (Engine API):} Quando un nodo remoto riceve il blocco, il suo CL estrae il payload di esecuzione e lo invia al proprio EL locale richiamando il metodo RPC \texttt{engine\_newPayload}.
    \item \textbf{Ri-esecuzione Nativa (EL):} L'EL deve verificare rigorosamente il blocco. Carica lo stato antecedente dal proprio disco, controlla la validità del base fee, la disponibilità di gas, i nonce, i saldi iniziali ed esegue sequenzialmente ogni transazione all'interno dell'EVM. Infine, ricalcola la radice dell'albero MPT.
    \item \textbf{Esito della Validazione:} Se il ricalcolo dello \textit{State Root} coincide con quello dichiarato dal Proposer, l'EL risponde al CL con lo status \texttt{VALID}.
    \item \textbf{Aggiornamento Fork Choice:} Avuta conferma della validità esecutiva, il CL invia il comando \texttt{engine\_forkchoiceUpdated}, istruendo l'EL a designare tale blocco come nuova \textit{head} della catena canonica.
\end{enumerate}

\section{Scalabilità: L2 vs.\ L1}

Storicamente, in risposta al severo limite intrinseco di transazioni al secondo dettato dal modello di ri-esecuzione $N$-of-$N$, la comunità di ricerca ha delegato il carico computazionale a reti secondarie. Questo approccio ha dato vita alla cosiddetta roadmap "Rollup-Centric". In questo paradigma, il Layer 1 (L1) sacrifica la scalabilità dell'esecuzione per mantenere un ruolo esclusivo come livello di massima sicurezza e garanzia della disponibilità dei dati. L'elaborazione massiva delle transazioni viene invece spostata \textit{off-chain} sui Layer 2 (L2), noti come \textit{Rollup}.

Esistono due approcci principali per garantire che l'esecuzione avvenuta su L2 sia valida quando viene finalizzata su L1:
\begin{itemize}
  \item \textbf{Optimistic Rollups:} Assumono per default che tutte le transazioni processate siano corrette. La sicurezza è garantita da un meccanismo di risoluzione delle dispute (\textit{Fraud Proofs}), che introduce tuttavia una latenza significativa (tipicamente di 7 giorni) prima che le transazioni raggiungano la finalità inoppugnabile sul L1.
  \item \textbf{ZK-Rollups:} Sfruttano le \textit{Validity Proofs} (prove crittografiche come SNARK o STARK) generate \textit{off-chain} per certificare l'integrità matematica di ogni batch di transazioni. Questo approccio garantisce la correttezza a priori e consente un \textit{settlement} sul L1 pressoché immediato senza periodi di contestazione.
\end{itemize}

\subsection{Criticità degli L2 attuali}

Nonostante il proliferare delle soluzioni L2, la loro architettura corrente presenta criticità sistemiche che contraddicono l'etica decentralizzata del L1. L'illusione di ereditare passivamente la sicurezza di Ethereum nasconde spesso l'impiego di \textit{Training Wheels}: meccanismi di salvaguardia centralizzati che conferiscono a pochi attori il potere di alterare il protocollo.

Analizzando il framework di classificazione \textit{L2Beat}\cite{l2beat}, si osserva che la quasi totalità degli L2 attuali si trova allo "Stage 0" o "Stage 1" di maturità decentralizzata. Questi network si affidano a \textit{Sequencer} operati da singole entità corporative e a \textit{Security Council} governati da chiavi \textit{multisig} capaci di sovrascrivere l'esecuzione o imporre aggiornamenti forzati degli smart contract di convalida.

Questa situazione ha portato a una profonda ricalibrazione teorica, evidenziata dalle recenti riflessioni di Vitalik Buterin: un L2 che non offre solide garanzie crittografiche \textit{trustless} non sta realmente scalando Ethereum, ma sta operando come un L1 indipendente mascherato da L2 tramite un bridge centralizzato. Per formalizzare questo concetto, è stato introdotto l'\textit{Ethereum Settlement Score} (ESS), una metrica che quantifica il grado in cui un L2 utilizza in modo inoppugnabile il livello di base per finalizzare le transazioni.

A livello di ecosistema, questa proliferazione orizzontale di Rollup imperfetti ha inoltre frammentato drasticamente la liquidità e distrutto la \textit{Synchronous Composability} — la capacità degli smart contract di interagire atomicamente tra loro all'interno del medesimo blocco, una proprietà fondamentale che aveva alimentato il successo della Finanza Decentralizzata (DeFi) sul Layer 1.

\subsection{Il Cambio di Paradigma: Scalare L1 tramite ZK}

Il superamento delle vulnerabilità degli L2 richiede un ritorno al potenziamento del Layer 1, abbandonando il paradigma in cui "tutti calcolano tutto" in favore di un modello in cui "uno calcola, tutti verificano". Questo rappresenta il salto architetturale dal calcolo ridondante ($N$-of-$N$) alla verifica crittografica ($1$-of-$N$). L'esecuzione di computazioni complesse richiede tempi lineari o polinomiali rispetto all'input ($O(N)$), mentre la verifica di una prova ZK richiede un tempo costante, riducendo l'onere computazionale asintotico a $O(1)$.

Questo principio è stato formalizzato dal ricercatore della Ethereum Foundation Justin Drake \cite{lean_ethereum} attraverso la dicotomia architetturale nota come \textbf{Beast Mode vs. Fort Mode}:
\begin{itemize}
  \item \textbf{Beast Mode:} Prevede l'esternalizzazione dell'esecuzione massiva e della generazione delle prove crittografiche a \textit{Provers} e \textit{Builders} operanti su infrastrutture di livello data-center. L'obiettivo prestazionale è spingere il limite di elaborazione del L1 verso la frontiera del \textit{GigaGas/s} (fino a 10.000 TPS), superando i limiti dei server general-purpose.
  \item \textbf{Fort Mode:} Rappresenta il baluardo della decentralizzazione. I validatori cessano di eseguire l'EVM nativamente per diventare validatori "ultra-leggeri". Grazie alla asimmetria computazionale fornita dalla crittografia ZK, questi nodi possono operare su hardware di consumo (laptop, smartphone o persino smartwatch) limitandosi a certificare matematicamente le prove generate dalle "Beast".
\end{itemize}

\vspace{4mm}

\begin{table}[H]
    \centering
    \renewcommand{\arraystretch}{1.4}
    \small
    \begin{tabular}{|p{0.25\textwidth}|p{0.32\textwidth}|p{0.32\textwidth}|}
        \hline
        \textbf{Caratteristica} & \textbf{Beast Mode} & \textbf{Fort Mode} (ZeroNethermind / Validators) \\
        \hline
        \textbf{Ruolo} & Esegue computazioni pesanti e genera prove crittografiche. & Mette in sicurezza la chain verificando le prove \\
        \hline
        \textbf{Hardware} & Specialized Hardware. & Laptop consumer, tablet, smartphone \\
        \hline
        \textbf{Storage (Stato L1)} & Terabyte & 0 GB (MemDB) \\
        \hline
        \textbf{Complessità Computazionale} & $O(N)$ - proporzionale al volume delle transazioni & $O(1)$ \\
        \hline
        \textbf{Obiettivo di Rete} & Scalare Ethereum a 10k+ TPS (GigaGas). & Resistenza incompromissibile alla censura e sicurezza \\
        \hline
    \end{tabular}
    \caption{Confronto architetturale tra il paradigma Beast Mode e Fort Mode}
    \label{tab:beast_vs_fort}
\end{table}

La transizione verso questa architettura è prevista in due fasi: una prima fase di \textit{optional proofs}\cite{eip8025}, in cui i nodi possono scegliere volontariamente di verificare i blocchi tramite prove ZK anziché ri-eseguirli, seguita da una fase di \textit{mandatory proofs}, in cui la verifica crittografica diviene il meccanismo predefinito e la ri-esecuzione un'opzione per ruoli specializzati (RPC, block building). ZeroNethermind anticipa la prima fase, dimostrando che un client EL può già oggi operare in modalità proof-only.

\section{Stateless clients}

Una delle proprietà più potenti di una rete blockchain è la capacità, per chiunque, di eseguire un nodo sul proprio computer e verificare in modo indipendente la correttezza della chain. Anche se la quasi totalità dei nodi in consenso dovesse accordarsi su regole malevole, un nodo in piena validazione rifiuterebbe categoricamente la chain compromessa. Tuttavia, affinché questa dinamica di sicurezza rimanga valida, l'esecuzione di un nodo deve rimanere accessibile a una massa critica di utenti.

Oggi, eseguire un nodo su hardware consumer è possibile, ma sempre più complesso a causa della mole di dati da archiviare. La fase della roadmap di Ethereum denominata "The Verge" si pone l'obiettivo di rivoluzionare questo aspetto, rendendo la validazione talmente leggera da poter essere eseguita di default su dispositivi mobili, wallet browser o perfino smartwatch.

\begin{figure}[H]
    \centering
    \includegraphics[width=1\textwidth]{img/verge.jpg}
    \caption{The Verge, 2023 roadmap \cite{verge_2024}}
    \label{fig:verge}
\end{figure}

Il problema fondamentale che The Verge cerca di risolvere è la crescita incontrollata dei dati di stato. Attualmente, un client Ethereum deve archiviare centinaia di gigabyte di dati per poter validare i blocchi, con un incremento di circa 30 GB di dati grezzi all'anno, a cui si aggiungono i metadati per aggiornare efficientemente il database. Questo attrito tecnologico allunga i tempi di sincronizzazione a giorni interi e disincentiva la partecipazione dei nodi indipendenti.

Originariamente, The Verge mirava a sostituire i Merkle Patricia Trie con i Verkle Trees, strutture dati che consentono prove molto più compatte. Oggi, la visione si è ampliata per includere la verifica dell'intera esecuzione di Ethereum tramite prove crittografiche SNARK, puntando a due obiettivi chiave:
\begin{itemize}
    \item \textbf{Client Stateless:} I client in piena validazione e i nodi di staking non dovranno richiedere più di una manciata di gigabyte di archiviazione.
    \item \textbf{Verifica su dispositivi ultra-leggeri:} Nel lungo termine, la convalida della catena consisterà unicamente nel download di una minima porzione di dati e nella verifica matematica di una SNARK.
\end{itemize}

\subsection{Stateless Ethereum}
"Stateless Ethereum" rappresenta un aggiornamento radicale del protocollo in cui i blocchi diventano unità di esecuzione self-contained. Il nodo non è più obbligato a scaricare l'intero stato della rete, poiché tutte le informazioni necessarie per l'esecuzione e la validazione vengono impacchettate direttamente all'interno del blocco stesso tramite prove crittografiche (i cosiddetti \textit{witnesses}).

I benefici principali di questa architettura si diramano su molteplici fronti:
\begin{itemize}
    \item \textbf{Scalabilità:} Rimuovendo il collo di bottiglia dell'I/O legato all'accesso al database su disco, i validatori possono processare molte più transazioni, permettendo un innalzamento sicuro del Gas Limit e quindi un TPS (Transaction Per Second) più elevato.
    \item \textbf{Decentralizzazione:} L'abbattimento dei requisiti hardware permette l'ingresso di nuovi attori nella sicurezza della rete, aprendo la strada a client ultra-leggeri \textit{trustless} e a dinamiche flessibili come il \textit{rainbow staking} o la creazione di pool di validazione privati.
    \item \textbf{Sincronizzazione Immediata:} Non necessitando del database storico, un nodo necessita soltanto del blocco corrente per unirsi attivamente al network.
\end{itemize}

\subsection{Relazione tra Verifica dei Blocchi e Stato}
Per comprendere appieno la relazione tra lo stato di Ethereum e la verifica dei blocchi, si consideri uno scenario in cui si debba validare un blocco contenente un'unica transazione: il trasferimento di 10 ETH dall'account A all'account B. Omettendo dettagli secondari, la logica di base impone di verificare che l'account A possieda un saldo sufficiente per coprire sia l'importo trasferito che le commissioni di rete (\textit{gas fees}). In caso contrario, la transazione deve fallire.

È fondamentale notare che, al momento della ricezione di un blocco, un client Execution Layer non può prevedere a priori quali account saranno coinvolti. Di conseguenza, in un'architettura tradizionale, il client è costretto a mantenere in memoria le informazioni relative a tutti gli account Ethereum esistenti. L'assenza di tali dati renderebbe impossibile la verifica del blocco per un client non equipaggiato, penalizzando in particolar modo le architetture \textit{stateless}. Le transazioni più complesse, come quelle che invocano \textit{smart contract}, seguono la medesima logica: i client devono disporre del bytecode dei contratti e del relativo stato di archiviazione (\textit{storage state}), poiché i blocchi possono teoricamente contenere transazioni che interagiscono con qualsiasi contratto distribuito sulla rete.

L'intuizione principale alla base del paradigma \textit{stateless} è che la porzione di stato effettivamente necessaria per verificare un singolo blocco è di ordini di grandezza inferiore rispetto all'intero \textit{World State}. Esistono porzioni dello stato di Ethereum che non vengono interrogate da anni, ma che devono comunque essere archiviate per l'eventualità che una nuova transazione possa richiederne l'accesso. Questa problematica fa emergere tangenzialmente la necessità di implementare meccanismi di \textit{State Expiry} (scadenza dello stato) a livello di protocollo.

\subsection{Validatori Stateless}
Affinché un nodo possa validare tradizionalmente un blocco, deve possedere lo stato antecedente (pre-state). In un modello fully-SNARKified, l'Execution Layer Client non deve eseguire alcun codice EVM: deve semplicemente validare una prova crittografica che attesta la corretta transizione dal \texttt{pre\_state\_root} al \texttt{post\_state\_root} applicando l'hash del blocco corrente.

Questa dinamica abilita la creazione dei \textbf{Validatori Stateless}. Poiché non dispongono del database completo, questi nodi non possono operare come \textit{Block Builders}, un compito che richiede l'accesso completo allo stato per simulare transazioni ed estrarre valore. Tuttavia, grazie a meccanismi già esistenti come la \textit{Proposer-Builder Separation} (PBS) o MEV-boost, un client stateless può delegare la costruzione a terze parti, ricevere il blocco candidato con la relativa prova, verificarne l'integrità in tempo costante e proporlo alla rete.

\subsection{Architettura del Client Stateless}
Il vantaggio ingegneristico di un client stateless risiede nella sua estrema semplicità architetturale rispetto a un full node. Passando al paradigma di pura verifica, interi moduli del client diventano obsoleti o superflui.

Osservando la struttura di un nodo stateless, notiamo una drastica riduzione delle componenti attive:
\begin{figure}[H]
    \centering
    \begin{tikzpicture}[
        box/.style={draw=black, very thick, dotted, minimum width=3cm, minimum height=2cm, align=center, font=\sffamily},
        req/.style={box, fill=green!30},
        opt/.style={box, fill=yellow!30},
        not/.style={box, fill=red!20}
    ]
    
    % --- Outer Border & Title ---
    \draw[thick] (-2, -6.5) rectangle (9, 2.5);
    \node[font=\sffamily\Large] at (3.5, 1.5) {Stateless architecture};

    % --- Row 1 ---
    \node[req] (api) at (0,0)      {engine api};
    \node[opt] (rpc) at (3.5,0)    {rpc};
    \node[not] (tx)  at (7,0)      {tx pool};

    % --- Row 2 ---
    \node[not] (evm)  at (0,-2.5)  {evm};
    \node[not] (tree) at (3.5,-2.5){tree};
    \node[not] (db)   at (7,-2.5)  {db};

    % --- Row 3 ---
    \node[not] (sync) at (0,-5)    {sync};
    \node[not] (p2p)  at (3.5,-5)  {p2p};
    \node[not] (blk)  at (7,-5)    {blockchain};

    % --- Legend ---
    \node[req, minimum width=1.5cm, minimum height=0.8cm] (l1) at (0,-7.5) {};
    \node[right, font=\sffamily] at (1,-7.5) {required};

    \node[opt, minimum width=1.5cm, minimum height=0.8cm] (l2) at (0,-8.7) {};
    \node[right, font=\sffamily] at (1,-8.7) {optional / partial};

    \node[not, minimum width=1.5cm, minimum height=0.8cm] (l3) at (0,-9.9) {};
    \node[right, font=\sffamily] at (1,-9.9) {not required};
    
    \end{tikzpicture}
    \caption{Stateless Client Architecture}
    \label{fig:stateless_arch}
\end{figure}
\begin{itemize}
    \item \textbf{Componenti eliminate (Not Required):} Il client viene snellito rimuovendo l'EVM (non essendovi esecuzione nativa), la Mempool delle transazioni (\textit{Tx Pool}), le logiche di aggiornamento dell'albero di stato (\textit{Tree}), il motore di persistenza su disco (\textit{DB}), la sincronizzazione storica (\textit{Sync}) e l'infrastruttura di propagazione Gossip (\textit{P2P} e \textit{Blockchain historical data}).
    \item \textbf{Componenti essenziali (Required):} L'interfaccia di comunicazione \textit{Engine API} rimane il cuore pulsante per interagire con il Consensus Layer.
    \item \textbf{Componenti opzionali:} Il servizio \textit{RPC} può essere mantenuto per interrogazioni limitate, sebbene non possa rispondere a query dipendenti dallo stato globale.
\end{itemize}

Per l'integrazione di questi client, stanno emergendo due pattern principali:
\begin{enumerate}
    \item \textbf{Separate Clients:} Il client stateless (CL+EL) si connette a peer \textit{stateful} tramite un sotto-protocollo di rete. L'esecuzione e la verifica sono disaccoppiate ma rimangono interattive.
    \item \textbf{Embedded Clients:} Un singolo client combina la logica di consenso (CL) con un Verificatore EL leggero incorporato (es. una zkVM o un esecutore stateless). Questo client riceve blocchi e prove tramite la rete gossip e li verifica interamente in locale, un modello che rispecchia l'architettura implementata in questa tesi.
\end{enumerate}

\section{Zero Knowledge}

L'evoluzione della crittografia informatica può essere suddivisa in due generazioni distinte. La "Prima Generazione" era caratterizzata da algoritmi \textit{special-purpose}, realizzati artigianalmente per assolvere compiti altamente specifici, come le funzioni di hash per l'integrità o le firme digitali per l'autenticazione delle transazioni. L'avvento della "Seconda Generazione" segna invece il passaggio alla crittografia \textit{general-purpose}. Introdotta formalmente alla fine degli anni Ottanta dai crittografi Goldwasser, Micali e Rackoff, la crittografia Zero-Knowledge (ZK) ne rappresenta la massima espressione, conferendo l'abilità inedita di formulare asserti crittografici su computazioni del tutto arbitrarie e Turing-complete. Una dimostrazione a conoscenza zero (Zero-Knowledge Proof, ZKP) permette a un'entità di provare matematicamente la veridicità di un'affermazione senza rivelare alcuna informazione sottostante o dato segreto. Questa tecnologia detiene il potere di trasferire il modello di fiducia dalle istituzioni centralizzate o dall'hardware specializzato direttamente all'infallibilità della matematica pura. 

A livello applicativo, i protocolli ZK abilitano una vasta gamma di casi d'uso trasversali. In ambito di identità digitale (ZK-ID e ZK-KYC), permettono a un utente di dimostrare il superamento di controlli anagrafici (ad esempio, l'essere maggiorenne) o il possesso di fondi sufficienti per una transazione, senza esporre la propria data di nascita o il saldo bancario reale. Nel settore finanziario istituzionale, le ZKP garantiscono la conformità alle normative antiriciclaggio (AML) e alle sanzioni, mantenendo contemporaneamente private le posizioni dei clienti e le strategie di trading. Nel panorama del calcolo distribuito in cloud, consentono inoltre di delegare carichi di lavoro computazionalmente onerosi a server non fidati, i quali restituiscono il risultato accompagnato da un certificato matematico di corretta esecuzione verificabile in tempo reale.  

È proprio su quest'ultimo aspetto che si basa la transizione di Ethereum verso la \textit{statelessness} e la scalabilità massiva, attraverso l'adozione di prove crittografiche succinte e non interattive (SNARG/SNARK). Il principio matematico fondante di questi protocolli in ottica di scalabilità risiede nell'introduzione di una netta asimmetria computazionale tra le entità partecipanti: da un lato figura il \textit{Prover}, una macchina dotata di altissima capacità di calcolo incaricata di eseguire la computazione e generare la prova crittografica; dall'altro opera il \textit{Verifier}, un nodo con risorse hardware limitate il cui unico onere è controllare rigorosamente la validità della prova ricevuta.  

In ambito blockchain, e in particolare nell'ecosistema Ethereum, è tuttavia doverosa una precisazione accademica formale: l'utilizzo del termine "Zero-Knowledge" (ZK) rappresenta spesso un \textit{misnomer} (un'inesattezza terminologica). La proprietà fondamentale che le cosiddette "zkEVM" sfruttano per abbattere il collo di bottiglia dell'esecuzione non è la \textit{privacy} (ovvero l'occultamento dei dati transazionali), bensì la \textbf{succintezza} (\textit{succinctness}). Di conseguenza, a livello ingegneristico, architetture come ZeroNethermind andrebbero più correttamente definite come \textit{Succinct EVMs} o \textit{SNARG-EVMs}.

\subsection{L'Interfaccia di una zkVM}

Dal punto di vista dell'integrazione di sistema, una zkVM opera come un costrutto crittografico che espone un'interfaccia standardizzata definita da tre funzioni principali:
\begin{itemize}
    \item \texttt{Prepare(f)} $\rightarrow$ \texttt{vk}: Durante la fase di setup, l'algoritmo riceve in input la descrizione formale del programma \texttt{f} (come la funzione di transizione di stato della EVM) e genera una chiave di verifica pubblica (\textit{verification key}, \texttt{vk}) di dimensioni estremamente ridotte.
    \item \texttt{Prove(f, x)} $\rightarrow$ \texttt{(y, proof)}: Funzione eseguita dal \textit{Prover}. Essa calcola l'output \texttt{y} a partire dallo stato di input \texttt{x} e produce contestualmente la prova crittografica (\texttt{proof}) che attesta la corretta esecuzione di \texttt{f(x)}.
    \item \texttt{Verify(vk, x, y, proof)} $\rightarrow$ \texttt{{0, 1}}: Funzione eseguita dal \textit{Verifier}. Presa in carico la chiave di verifica, gli input/output pubblici e la prova, l'algoritmo restituisce un valore booleano che certifica o rigetta la validità della transizione in un tempo sublineare.
\end{itemize}

\subsection{Le Tre Proprietà Fondamentali}

Affinché la rete possa delegare in sicurezza l'esecuzione \textit{trustless}, il sistema di proving sottostante deve garantire tre proprietà matematiche inderogabili:
\begin{enumerate}
    \item \textbf{Completeness (Completezza):} Se l'affermazione computazionale è vera e il \textit{Prover} si comporta in modo onesto, il \textit{Verifier} accetterà sempre la prova. Questa proprietà è cruciale per preservare la \textit{liveness} della blockchain, evitando falsi negativi che bloccherebbero il progresso della catena.
    \item \textbf{Soundness (Solidità/Robustezza):} Se l'affermazione è falsa, nessun avversario computazionalmente limitato può generare una prova valida, se non con probabilità trascurabile. Nelle zkVM, si fa ricorso alla \textit{Knowledge Soundness}, la quale impone che il \textit{Prover} non solo calcoli il risultato, ma conosca effettivamente l'intera traccia dell'esecuzione (\textit{witness}) per poter costruire la prova.
    \item \textbf{Succinctness (Succintezza):} Rappresenta la chiave di volta per risolvere il trilemma della scalabilità. La dimensione della prova e, soprattutto, il tempo di verifica computazionale risultano sublineari rispetto al tempo necessario per la ri-esecuzione nativa $O(N)$ del programma, assestandosi su tempi di risoluzione costanti ($O(1)$) o polilogaritmici.\end{enumerate}

\subsection{Problemi Architetturali Risolti dalle zkEVM}

La scalabilità di qualsiasi architettura blockchain si scontra universalmente con tre vincoli fondamentali: \textit{Compute}, \textit{I/O} e \textit{Bandwidth}. L'introduzione delle zkEVM interviene per mitigare congiuntamente queste limitazioni architetturali. In primo luogo, il vincolo del calcolo viene risolto direttamente dalla natura stessa delle prove ZK, che comprimono l'onere computazionale in una verifica crittografica pressoché istantanea. In secondo luogo, il vincolo dell'I/O viene risolto indirettamente: non dovendo ri-eseguire le transazioni nell'EVM, i nodi sono sollevati dall'obbligo di caricare e aggiornare costantemente l'enorme database di stato locale, abilitando di fatto operazioni parzialmente o totalmente \textit{stateless}. Infine, il problema della larghezza di banda viene affrontato accoppiando il paradigma ZK con meccanismi di Data Availability Sampling (DAS), permettendo alla rete di attestare la disponibilità dei dati e dei blocchi compressi senza forzare i nodi al download integrale di tutte le transazioni storiche.

Nel protocollo attuale, il throughput è strettamente vincolato dal \textit{Block Gas Limit}. Un suo incremento richiederebbe ai validatori hardware sempre più costoso per sostenere la ri-esecuzione completa, accentrando il consenso. Nelle zkEVM, invece, un singolo \textit{builder} esegue l'intero blocco e un \textit{prover} ne genera la prova crittografica. I validatori si limitano al \textit{proof checking}, un'operazione asintoticamente più economica. Questo paradigma $1$-of-$N$ permette di innalzare il limite di gas in totale sicurezza: espandendo la capacità dei blocchi si spalma l'alta domanda su una maggiore disponibilità di spazio, stabilizzando e abbattendo i prezzi delle commissioni.

Parallelamente, delegare il carico computazionale ai \textit{prover} previene i ritardi nella \textit{finality}. Poiché la verifica crittografica richiede un tempo costante e indipendente dalla complessità del blocco, i validatori non rischiano più di perdere gli slot a causa di ri-esecuzioni troppo lente. La barriera hardware viene così abbattuta radicalmente, permettendo persino a dispositivi modesti (come un Raspberry Pi) di partecipare alla messa in sicurezza della rete e preservandone l'assoluta diversità.

Separando nettamente la fase di esecuzione massiva da quella di validazione, le zkEVM trasformano radicalmente il modello di scalabilità di Ethereum verso un'architettura unificata. Questa dicotomia genera un mercato dell'esecuzione altamente competitivo che incentiva l'ottimizzazione dell'hardware specializzato per i \textit{prover}, garantendo parallelamente una massiccia diversità dei validatori tramite l'empowerment di nodi leggeri. Il risultato finale è un throughput elastico e sicuro, che consente alla capacità del Layer 1 di crescere senza mai compromettere la decentralizzazione e l'accessibilità della rete.

\section{Lavori Correlati}

L'architettura e i principi di design alla base del progetto presentano radici in un filone di ricerca ampio e in rapida evoluzione, che spazia dalla teoria dei sistemi distribuiti fino alle più recenti proposte di aggiornamento del protocollo Ethereum (EIP). Questa sezione analizza la letteratura e i progetti fondamentali che fanno da preludio all'implementazione di un nodo validatore Layer 1 (L1) basato su zkEVM, inserendo il prototipo \textit{ZeroNethermind} nel contesto della roadmap ufficiale.

\subsection{Il Trilemma della Blockchain e i Client Leggeri}

Storicamente, il design dei registri distribuiti è stato dominato dal concetto del \textit{Blockchain Trilemma}, un'euristica che postula l'impossibilità di massimizzare simultaneamente scalabilità, sicurezza e decentralizzazione. Tuttavia, recenti analisi formali \cite{trilemma} hanno confutato la natura assoluta di questo trilemma, dimostrando che esso deriva spesso da un'equivocazione semantica e dall'errata applicazione di teoremi propri dei sistemi distribuiti classici a protocolli basati su prove crittografiche. In reti come Ethereum, la sicurezza non emerge dalla ridondanza simmetrica della computazione di tutti i nodi, ma dalla disseminazione di prove crittografiche. Di conseguenza, la scalabilità non è un compromesso obbligato contro la decentralizzazione, ma un risultato ingegneristico raggiungibile disaccoppiando la produzione delle prove dalla loro verifica.

Questa necessità di disaccoppiamento richiama il concetto storico di \textit{Light Client}. Come sistematizzato nel documento \textit{SoK: Blockchain Light Clients} \cite{blockchain_light_clients}, i client leggeri permettono di interagire con la blockchain senza memorizzarne l'intera cronologia. Mentre i light client tradizionali si affidano ad assunzioni di fiducia, la ricerca moderna si sta orientando sull'utilizzo di accumulatori crittografici e zk-SNARKs per permettere a dispositivi con risorse limitate di verificare matematicamente l'integrità del sistema in modalità pienamente \textit{trustless}.

\subsection{The Verge e la "Statelessness" di Ethereum}

Il superamento dei limiti computazionali dei nodi è il nucleo della fase di sviluppo della roadmap di Ethereum nota come \textit{The Verge} \cite{verge_2024}. L'obiettivo primario è democratizzare la validazione, rendendola così efficiente da poter essere eseguita su dispositivi di largo consumo (smartphone o smartwatch). 

Attualmente, un nodo Ethereum richiede la memorizzazione di un database di stato in continua crescita, creando un collo di bottiglia I/O che innalza le barriere all'ingresso e spinge verso la centralizzazione. Il paradigma \textit{Stateless Ethereum} \cite{stateless_book} risolve questo problema trasformando i blocchi in unità di esecuzione auto-contenute: anziché interrogare un database locale, il nodo stateless riceve un blocco accompagnato da un \textit{witness} (testimone) e da una prova crittografica che ne attesta la validità. 

A supporto di questa profonda transizione strutturale, è stata proposta l'\textbf{EIP-4762} (\textit{Statelessness gas cost changes}) \cite{eip4762}. Questa proposta introduce una rimodulazione dei costi del gas per riflettere accuratamente l'onere crittografico della generazione dei witness, disincentivando operazioni di stato che renderebbero le prove crittografiche troppo pesanti o complesse da generare.

\subsection{Il Panorama dei Layer 2 e l'Analisi dei Rischi}

Fino ad oggi, l'onere della scalabilità è stato delegato ai Layer 2 (L2) tramite la \textit{rollup-centric roadmap}. Piattaforme analitiche come \textbf{L2BEAT} \cite{l2beat} svolgono un ruolo cruciale nel valutare i compromessi di sicurezza di queste architetture, dimostrando come i rollup attuali si collochino su uno "spettro di fiducia". Molte soluzioni operano ancora in \textit{Stage 1}, basandosi su sequencer centralizzati, bridge controllati da multisig e \textit{Security Councils} autorizzati a forzare l'aggiornamento dello stato. 

Per mitigare questi vettori di centralizzazione e i rischi derivanti da un ecosistema frammentato, la ricerca sta virando verso l'integrazione nativa delle prove ZK nel Layer 1 (\textit{Enshrined zkEVM}). Questa mossa mira a unificare le garanzie di sicurezza dell'ecosistema, dimostrando che la tecnologia zkVM, maturata sui Layer 2, è pronta per il livello base del protocollo.

\subsection{Lean Ethereum: Il Cambio di Paradigma}

Il cambio di paradigma di Lean Ethereuem introdotto da Justin Drake, che passa dal modello \textit{"Trust but Verify"} al modello \textit{"Verify the Math"}, è delineato nella \textbf{L1-zkEVM Roadmap 2026} \cite{l1_zkevm_roadmap}. Tale documento definisce il flusso end-to-end che parte da un \textit{Execution Client} (che genera l'\textit{ExecutionWitness}), passa per un programma ospite zkVM e un Prover, per finire al nodo validatore che ne asserisce la validità.

In questa direzione, l'\textbf{EIP-8025} (\textit{Optional Execution Proofs}) \cite{eip8025} rappresenta un passo fondamentale: introduce una modifica opzionale a livello di consenso che permette ai nodi Beacon di verificare la validità dei payload di esecuzione tramite prove ZK, senza dover interrogare un client di esecuzione locale.

\subsection{L'Ecosistema zkVM: Implementazioni e Competizione}

L'implementazione pratica del \textit{real-time proving} poggia su macchine virtuali a conoscenza zero (zkVM) capaci di interpretare set di istruzioni, monitorate attivamente da portali come \textit{zkevm.fyi} \cite{zkevm_book}. Tra i progetti di punta figurano:
\begin{itemize}
    \item \textbf{SP1}\cite{sp1} (Succinct Labs): Una zkVM \textit{general-purpose} ad alte prestazioni basata sull'architettura RISC-V, che permette di generare prove per programmi compilati in Rust (o tramite LLVM).
    \item \textbf{ZisK}\cite{zisk} (Polygon): Un'architettura zkVM altamente specializzata, progettata specificamente per l'esecuzione rapida e la generazione di prove per transazioni compatibili con l'EVM.
\end{itemize}
La coesistenza e la competizione tra questi \textit{prover} sono vitali per garantire la \textit{client diversity} anche a livello crittografico, scongiurando singoli punti di fallimento (es. bug in un circuito ZK) richiedendo \textit{multi-proofs} per l'accettazione di un blocco.

\subsection{Prototipi e L1 Strawmap}

La validità di questo approccio è stata dimostrata sul campo durante conferenze di settore come il Devconnect. In tale occasione, è stata presentata \textbf{zkLighthouse} \cite{zkLighthouse_demo}, una versione modificata del client di consenso Lighthouse. Operando su hardware di largo consumo, questo nodo ha bypassato completamente l'uso di un \textit{Execution Client} classico (come Geth o Nethermind), validando i blocchi della Mainnet in tempo reale. Il sistema si limitava ad ascoltare un canale P2P dedicato, verificando in millisecondi le prove prodotte da cluster di prover esterni.

L'intero sforzo architetturale si condensa infine nella \textbf{L1 Strawmap} (\textit{strawmap.org}) \cite{strawmap}, una bozza di roadmap mantenuta dal team della Ethereum Foundation. Questa mappa unifica gli aggiornamenti in tre filoni (Consenso, Dati, Esecuzione) e fissa le pietre miliari del protocollo, tra cui la transizione a un \textit{Gigagas L1} (10.000 TPS) resa possibile proprio dall'integrazione delle prove zkEVM e del \textit{real-time proving}.

Lo sviluppo del presente lavoro di tesi si inserisce esattamente all'incrocio di queste tecnologie, proponendo la re-ingegnerizzazione di un \textit{Execution Layer} \textit{stateless} in grado di agire come puro validatore di prove all'interno del framework \textit{Fort Mode} delineato dal futuro del protocollo Ethereum.

% ============================================================
%  CAPITOLO 3: ZERONETHERMIND
% ============================================================
\chapter{ZeroNethermind}\label{chap:zeronethermind}

Questo capitolo presenta la progettazione ad alto livello di ZeroNethermind, sviluppato per dimostrare la fattibilità di un Execution Layer completamente \textit{stateless} basato sulla verifica di prove crittografiche Zero-Knowledge. L'analisi parte dall'esame del client \texttt{Nethermind} e dell'infrastruttura esterna \texttt{ethproofs.org}, per poi descrivere il modello architetturale del "Validatore Stateless" e la progettazione dei componenti core del sistema. I dettagli implementativi saranno trattati nel \texttt{capitolo \ref{chap:implementazione}}, mentre la validazione sperimentale nel \texttt{capitolo \ref{chap:validazione}}.

\section{Nethermind Client}

Nethermind è uno dei principali client dell'Execution Layer per Ethereum, sviluppato interamente in C\#/.NET. Il progetto è entrato in produzione nel 2019 e rappresenta, insieme a Geth (Go Ethereum) e Reth (Rust Ethereum), uno dei tre client EL più diffusi nella rete mainnet, contribuendo attivamente alla diversità del protocollo.

\subsection{Requisiti di Sistema}

I requisiti hardware\cite{nethermind_requirements} per eseguire un nodo Nethermind in piena validazione sono significativi e in costante crescita:
\begin{itemize}
    \item \textbf{CPU:} 4 o più core.
    \item \textbf{RAM:} Almeno 16\,GB (32\,GB consigliati per prestazioni ottimali).
    \item \textbf{Storage:} Oltre \textbf{2\,TB di NVMe SSD} ad alte prestazioni (min 10.000 IOPS).
\end{itemize}

Nethermind impiega \textbf{RocksDB}\cite{nethermind_database}, un motore di database chiave-valore ad alte prestazioni, per la persistenza su disco. I dati sono organizzati in circa sei categorie (State, Receipts, Blocks, Headers, Code, Metadata), con il solo database di stato che supera i 150\,GB. L'intera base dati, comprensiva dello storico dei blocchi e delle ricevute, raggiunge e supera 1\,TB. Questa massiccia impronta di storage, unita alla necessità di IOPS elevati, rappresenta la barriera primaria che rischia di escludere l'hardware di livello consumer dalla partecipazione attiva alla rete in futuro.

\subsection{Il Flusso di Validazione Standard}

Come descritto nel \texttt{capitolo \ref{chap:introduzione}}, la comunicazione tra Consensus Layer (CL) ed Execution Layer (EL) avviene tramite l'\textbf{Engine API}, un insieme di metodi JSON-RPC su canale autenticato. Le due chiamate centrali nel flusso di validazione sono:
\begin{enumerate}
    \item \texttt{engine\_newPayload}: il CL invia un \texttt{ExecutionPayload} contenente i dati del blocco. L'EL standard \textbf{ri-esegue tutte le transazioni} nell'EVM, ricalcola lo \texttt{stateRoot} e restituisce uno status \texttt{VALID} o \texttt{INVALID}.
    \item \texttt{engine\_forkchoiceUpdated}: il CL comunica all'EL gli hash dei blocchi \texttt{head}, \texttt{safe} e \texttt{finalized}. L'EL aggiorna di conseguenza la propria catena canonica.
\end{enumerate}

Il costo computazionale di questo flusso è dominato dalla prima chiamata: la ri-esecuzione nativa impone una complessità $O(N_{gas})$ per ogni blocco, dove $N_{gas}$ è la quantità totale di gas consumata. Questo costo è prevalentemente I/O-bound, essendo determinato dalle continue operazioni di lettura e scrittura sul database per accedere allo stato pre-esistente e persistere il nuovo stato risultante.

% TODO: spostare in Stato dell'Arte
\section{Ethproofs.org}

Ethproofs.org\cite{ethproofs} è un \textit{Block Proof Explorer} per Ethereum: una piattaforma pubblica che aggrega prove crittografiche ZK generate da molteplici team di proving indipendenti per i blocchi della Mainnet. Lanciato da parte dell'Ethereum Foundation, ethproofs.org svolge un ruolo analogo a quello di un block explorer tradizionale, ma specializzato nel tracciamento delle ZK proofs.

La piattaforma oltre a presentare un sito web interattivo, espone un'API REST pubblica che consente di interrogare i metadati delle prove disponibili per ciascun blocco e di scaricare sia le prove binarie sia le relative chiavi di verifica (\textit{Verification Keys}, VK).

\subsection{L'Ecosistema dei Prover}

Una delle caratteristiche più significative di ethproofs.org è la sua natura di \textbf{aggregatore multi-prover}. La piattaforma raccoglie prove generate da diversi cluster di proving, ciascuno basato su una zkVM (\textit{Zero-Knowledge Virtual Machine}) distinta. La tabella \ref{tab:provers} riassume i principali prover attualmente operativi sulla piattaforma.

\begin{table}[H]
\centering
\begin{tabular}{l l l}
\hline
\textbf{Prover} & \textbf{Famiglia zkVM} & \textbf{Sistema di Prova} \\
\hline
Zisk & Polygon Hermez & pil2-proofman \\
OpenVM & Axiom & STARK (verify-stark) \\
Pico & Pico Prism & KoalaBear Poseidon2 \\
SP1 Hypercube & Succinct Labs & Compressed STARK \\
Airbender & Matter Labs (zkSync) & Mersenne31 + Unified Recursion \\
\hline
\end{tabular}
\caption{Principali prover operativi su ethproofs.org.}
\label{tab:provers}
\end{table}

Questa client diversity è intenzionale. Molteplici team competono per generare prove valide nel minor tempo possibile, prefigurando la nascita di un vero e proprio \textit{Prover Market}\cite{prover_market}. 

\subsection{Le API e il Modello Dati}

L'interazione programmatica con ethproofs.org avviene attraverso tre endpoint REST principali:
\begin{enumerate}
    \item \texttt{GET /api/blocks/\{blockId\}/proofs}: restituisce i metadati delle prove disponibili per un dato blocco.
    \item \texttt{GET /api/proofs/download/\{proofId\}}: restituisce i file binari della prova crittografica.
    \item \texttt{GET /api/verification-keys/download/\{proofId\}}: restituisce i file binari della chiave di verifica associata.
\end{enumerate}

\subsection{Ruolo nel PoC}

Nell'architettura di ZeroNethermind, ethproofs.org riveste il ruolo di \textit{single source of truth} per le ZK proofs, eliminando così la necessità per il PoC di operare una propria infrastruttura di proving, la quale richiederebbe cluster GPU di alto profilo e competenze specialistiche nella generazione delle prove. ZeroNethermind si concentra esclusivamente sulla verifica, delegando interamente la generazione all'ecosistema esterno.

\section{Design e Componenti}

Il concetto fondante di ZeroNethermind è sintetizzato nel paradigma \textbf{``Double Zero''}, che combina due azzeramenti architetturali complementari:
\begin{itemize}
    \item \textbf{Zero State}: il client non mantiene alcun database di stato locale. Operando in modalità \texttt{MemDb} (memoria volatile), tutta la persistenza è effimera e limitata ai soli header dei blocchi validati. L'impronta su disco è pertanto prossima allo zero.
    \item \textbf{Zero Knowledge}: la validità dei blocchi non è asserita tramite ri-esecuzione delle transazioni nell'EVM, bensì tramite la verifica crittografica di prove ZK-SNARK/STARK generate da prover esterni.
\end{itemize}

Lo spostamento di complessità è radicale: dalla ri-esecuzione con costo $O(N_{gas})$, proporzionale al volume di gas consumato nel blocco, alla verifica crittografica con costo $O(1)$, indipendente dalla complessità del blocco. Un blocco contenente una singola transazione richiede lo stesso tempo di verifica di un blocco contenente migliaia di transazioni.

\subsection{Architettura}

L'architettura di ZeroNethermind è organizzata in \textbf{tre strati} con responsabilità distinte, ciascuno isolato e comunicante attraverso interfacce ben definite. Questa stratificazione riflette sia le scelte ingegneristiche sia le necessità tecnologiche (interoperabilità cross-language tra C\# e Rust).

\begin{figure}[H]
\centering
\resizebox{\textwidth}{!}{%
\begin{tikzpicture}[
    >=Stealth,
    node distance=1.8cm and 2.5cm,
    font=\sffamily,
    % Stili base per i nodi
    core/.style={rectangle, draw=cyan!60!black, rounded corners=8pt, fill=cyan!10, thick, minimum width=3.8cm, minimum height=1.2cm, align=center},
    ext/.style={rectangle, draw=gray!80, rounded corners=8pt, fill=gray!10, thick, minimum width=3.2cm, minimum height=1.2cm, align=center},
    evm/.style={rectangle, draw=gray!40, rounded corners=8pt, fill=gray!5, text=gray!50, thick, minimum width=3.2cm, minimum height=1.2cm, align=center},
    db/.style={cylinder, draw=gray!70, shape border rotate=90, aspect=0.25, fill=gray!20, thick, minimum width=2.4cm, minimum height=1.8cm, align=center},
    rust/.style={rectangle, draw=orange!80!black, rounded corners=8pt, fill=orange!10, thick, minimum width=3.8cm, minimum height=1.2cm, align=center},
    crypto/.style={rectangle, draw=green!60!black, rounded corners=6pt, fill=green!10, thick, minimum width=2cm, minimum height=0.8cm, align=center},
    % Stili per gli strati (bounding box)
    layerbox/.style={draw=gray!50, dashed, thick, rounded corners=12pt, inner sep=16pt, fill=gray!5},
    % Stili per le frecce
    arrow/.style={->, thick, draw=blue!60!black},
    arrowdashed/.style={->, thick, dashed, draw=gray!50},
    zerocopy/.style={->, very thick, draw=orange!80!black}
]

% ================= COMPONENTI ESTERNI =================
\node[ext] (consensus) {Consensus Layer};
\node[ext, right=4.5cm of consensus] (ethproofs) {ethproofs.org};

% ================= STRATO 1: INTERCEPTOR =================
\node[core, below=2cm of consensus] (handler) {ZkNewPayloadHandler};
\node[evm, right=1.5cm of handler] (stdevm) {Standard EVM};

\draw[red!60, very thick] (stdevm.north west) -- (stdevm.south east);
\draw[red!60, very thick] (stdevm.south west) -- (stdevm.north east);

% ================= STRATO 2: VERIFICATORE =================
\node[core, below=2.5cm of handler] (verifier) {ZkValidationService};
\node[db, left=2.5cm of verifier] (memdb) {MemDb\\\textit{\small Zero-State}};

% ================= STRATO 3: BRIDGE FFI =================
\node[rust, below=3.2cm of verifier] (native) {native-zk-verifier\\(Rust)};

\node[crypto, below=1.5cm of native] (zisk) {Zisk};
\node[crypto, left=0.2cm of zisk] (openvm) {OpenVM};
\node[crypto, left=0.2cm of openvm] (pico) {Pico};
\node[crypto, right=0.2cm of zisk] (sp1) {SP1};
\node[crypto, right=0.2cm of sp1] (airbender) {Airbender};

% ================= BACKGROUND LAYERS =================
\begin{scope}[on background layer]
    \node[layerbox, fit=(handler) (stdevm), label={[font=\bfseries, text=black!60]above left:STRATO 1: Interceptor}] (layer1) {};
    \node[layerbox, fit=(memdb) (verifier), label={[font=\bfseries, text=black!60]above left:STRATO 2: Verifier}] (layer2) {};
    \node[layerbox, fit=(native) (pico) (airbender), label={[font=\bfseries, text=black!60]above left:STRATO 3: FFI Bridge}] (layer3) {};
\end{scope}

% ================= COLLEGAMENTI (FLUSSI) =================

\draw[arrow] (consensus) -- node[pos=0.2, left=0.2cm, font=\small, align=right] {Engine API} (handler);
\draw[arrowdashed] (handler) -- node[above=0.1cm, font=\small, text=gray!60] {Bypass} (stdevm);
\draw[arrow] (handler) -- node[left=0.2cm, font=\small] {Payload \& Context} (verifier);
\draw[arrow, <->] (verifier) -- node[above=0.1cm, font=\small] {State Queries} (memdb);
\draw[arrow] (verifier.east) -| node[pos=0.75, right=0.1cm, font=\small, align=left] {Fetching\\HTTP GET} (ethproofs.south);
\draw[zerocopy] (verifier) -- node[pos=0.55, right=0.2cm, font=\small, align=left] {Puntatori Unsafe\\(Zero-Copy FFI)} (native);
\draw[arrow, draw=green!60!black] (native.south) -- ++(0,-0.5) -| (zisk.north);
\draw[arrow, draw=green!60!black] (native.south) -- ++(0,-0.5) -| (openvm.north);
\draw[arrow, draw=green!60!black] (native.south) -- ++(0,-0.5) -| (pico.north);
\draw[arrow, draw=green!60!black] (native.south) -- ++(0,-0.5) -| (sp1.north);
\draw[arrow, draw=green!60!black] (native.south) -- ++(0,-0.5) -| (airbender.north);

\end{tikzpicture}%
}
\caption{Architettura core di ZeroNethermind}
\label{fig:zeronethermind_arch}
\end{figure}

La Figura \ref{fig:zeronethermind_arch} illustra l'architettura logica e il flusso di esecuzione del sistema, scomposto in tre macro-livelli:

\begin{itemize}
    \item \textbf{Strato 1 (Interceptor):} Rappresenta il punto di contatto con l'Engine API. Il gestore custom (\texttt{ZkNewPayloadHandler}) intercetta i payload in ingresso dal Consensus Layer, bypassando integralmente la tradizionale macchina a stati (\texttt{Standard EVM}) e azzerando l'onere computazionale nativo.
    \item \textbf{Strato 2 (Verifier):} Costituisce il cuore logico scritto in \texttt{C\#}. Il servizio orchestratore coordina il recupero parallelo delle prove tramite chiamate HTTP all'oracolo esterno (\texttt{ethproofs.org}). Inoltre, gestisce la persistenza effimera della catena: gli \textit{header} validati vengono salvati nel \texttt{MemDb}, permettendo al client di operare in uno stato di totale dipendenza dalla memoria RAM (\textit{Zero-State}), senza ricorrere ad alcun database su disco.
    \item \textbf{Strato 3 (FFI Bridge):} Rappresenta il confine di interoperabilità ad alte prestazioni tra l'ambiente \textit{managed} di .NET e le librerie crittografiche native. Per evitare colli di bottiglia causati dalla serializzazione e dal Garbage Collector, i payload binari (prove e chiavi di verifica) vengono passati allo strato \texttt{Rust} tramite puntatori bloccati in memoria (\textit{Zero-Copy FFI}). Questo strato agisce come dispatcher verso gli algoritmi di verifica specifici delle singole zkVM (Zisk, OpenVM, Pico, SP1, Airbender).
\end{itemize}

La Tabella \ref{tab:confronto_nethermind} sintetizza le differenze architetturali tra un nodo Nethermind standard e ZeroNethermind. Il confronto evidenzia come il PoC ribalti il profilo delle risorse richieste: dallo storage massivo e dall'I/O intensivo di un full node, si passa a un modello dominato dalla latenza di rete.

\begin{table}[H]
\centering
\begin{tabular}{l l l}
\hline
\textbf{Caratteristica} & \textbf{Nethermind Standard} & \textbf{ZeroNethermind} \\
\hline
Input & Blocco + Stato Completo & Blocco + Prova ZK \\
Azione & Ri-esecuzione Tx (EVM) & Verifica Locale \\
Storage di Stato & $\sim$2\,TB (RocksDB) & $\sim$0\,GB (MemDb, solo header) \\
Risorse & Alta CPU, Alto I/O SSD & Bassa CPU, Alta Rete \\
Bottleneck & Velocità Disco (IOPS) & Disponibilità Prove (Latenza) \\
Cold Start & Ore (Snap Sync) & Secondi \\
\hline
\end{tabular}
\caption{Confronto architetturale tra Nethermind Standard e ZeroNethermind.}
\label{tab:confronto_nethermind}
\end{table}

\subsection{Il Flusso di Dati Modificato}

Il flusso di validazione in ZeroNethermind differisce radicalmente da quello standard descritto nella sezione \ref{sec:validazione_standard}. Il processo segue questi step architetturali:

\begin{figure}[H]
    \centering
    \includegraphics[width=\textwidth]{img/zn_validation_flow.png}
    \caption{Flusso di Validazione di ZeroNethermind}
\end{figure}

\begin{enumerate}
    \item \textbf{Ricezione del blocco:} il CL riceve un blocco dalla rete gossip ed invoca \texttt{engine\_newPayload} verso ZeroNethermind, trasmettendo l'\texttt{ExecutionPayload}.
    \item \textbf{Pre-processing minimale:} ZeroNethermind decodifica il blocco e ne valida l'hash dell'header. A differenza del flusso standard, non vengono eseguiti i controlli che dipendono dal database di stato locale (verifica parent header, controllo pre-pivot, etc.).
    \item \textbf{Recupero delle prove:} il sistema interroga l'API di ethproofs.org per ottenere i metadati delle prove disponibili per il numero di blocco corrente.
    \item \textbf{Verifica crittografica:} per ciascuna prova disponibile, ZeroNethermind scarica i byte binari della prova e della chiave di verifica, quindi invoca il verificatore nativo (Rust) tramite FFI.
    \item \textbf{Consenso multi-prover:} il blocco è considerato valido se almeno il 50\% delle prove verificate risulta valido. Questo approccio mitiga il rischio di bug in un singolo circuito zkVM.
    \item \textbf{Risposta al CL:} ZeroNethermind restituisce lo status \texttt{VALID}, \texttt{INVALID} o \texttt{SYNCING} al Consensus Layer.
    \item \textbf{Aggiornamento della catena:} alla ricezione di \newline \texttt{engine\_forkchoiceUpdated}, il sistema aggiorna la testa canonica.
\end{enumerate}

È importante sottolineare che il nodo opera con un profilo di configurazione dedicato (\texttt{zk-mainnet}) che disabilita esplicitamente le componenti non necessarie — P2P, sincronizzazione, TxPool — e attiva la modalità MemDb. I dettagli di questa ottimizzazione saranno approfonditi nella sezione \ref{sec:config}.

Analogamente, la gestione del protocollo \texttt{SYNCING}, che regola l'interazione con il CL quando le prove non sono ancora disponibili e il nodo entra in modalità ``optimistic'', sarà descritta nella sezione \ref{sec:syncing}.


% ============================================================
%  CAPITOLO 4: IMPLEMENTAZIONE
% ============================================================
\chapter{Implementazione}\label{chap:implementazione}

Questo capitolo documenta il percorso ingegneristico dalla prima sperimentazione alla realizzazione del prototipo ZeroNethermind funzionante. L'implementazione è stata condotta in tre fasi incrementali, ciascuna costruita sulle fondamenta della precedente: la creazione di un verificatore ZK standalone (sezione \ref{sec:foundation}), la sua integrazione nel client Nethermind (sezione \ref{sec:integration}), e l'orchestrazione del servizio di validazione completo (sezione \ref{sec:validation_service}). Il capitolo si conclude con l'analisi del profilo di configurazione che trasforma il client in un validatore stateless (sezione \ref{sec:config}).

\section{Fase 1: The Foundation}\label{sec:foundation}

L'obiettivo della prima fase è stato costruire un verificatore ZK multi-prover standalone, denominato \texttt{EthProofValidator}, per dimostrare la fattibilità dell'interoperabilità .NET $\leftrightarrow$ Rust, indipendentemente dal client Nethermind.

\subsection{L'Ecosistema dei Verifier ZK}

Come descritto nel Capitolo~\ref{chap:introduzione}, l'ecosistema delle zkEVM è composto da molteplici team indipendenti, ciascuno dei quali ha sviluppato la propria zkVM e il relativo sistema di proving. Per ZeroNethermind, questi verifier sono trattati come \textit{black-box}: il PoC si interfaccia esclusivamente con le loro API di verifica, senza intervenire sulla logica interna.

Un aspetto tecnico rilevante è che tutti i principali framework di proving sono implementati in \textbf{Rust}, il linguaggio dominante per la computazione crittografica ad alte prestazioni. Questo crea una sfida d'integrazione significativa con il client Nethermind, che è interamente basato su C\#/.NET: è necessario un bridge affidabile ed efficiente tra i due runtime.

Le cinque zkVM attualmente integrate nel PoC sono: \textbf{Zisk} (Polygon Hermez), \textbf{OpenVM} (OpenVM.org), \textbf{Pico} (Brevis), \textbf{SP1 Hypercube} (Succinct) e \textbf{Airbender} (zkSync/Matter Labs). Ciascun team espone un'interfaccia di verifica concettualmente analoga — una funzione che accetta una prova e una chiave di verifica, e ritorna un risultato booleano — ma con differenze nei formati di serializzazione e nelle dipendenze interne. Queste differenze rendono necessaria la creazione di un \texttt{crate} Rust unificato (\texttt{native-zk-verifier}) che funge da adattatore comune.

\subsection{WASM vs Native FFI}\label{sec:wasm_vs_ffi}

Per realizzare il bridge C\# $\leftrightarrow$ Rust, sono state valutate due strategie architetturali: il caricamento dei verifier come moduli \textbf{WebAssembly} (WASM), come fatto da ethproofs.org\cite{ethproofs} per l'\textit{in-browser} validation, tramite il runtime \texttt{Wasmtime.NET}, e l'invocazione diretta tramite \textbf{Foreign Function Interface} (FFI) nativa con \texttt{P/Invoke}. La \ref{tab:wasm_vs_ffi} sintetizza il confronto.

\begin{table}[H]
\centering
\begin{tabular}{l p{5cm} p{5cm}}
\hline
\textbf{Aspetto} & \textbf{WASM} & \textbf{Native FFI (P/Invoke)} \\
\hline
Prestazioni & Buone (VM overhead). & \textbf{Massime} (esecuzione diretta su CPU). \\
Portabilità & Cross-platform nativo. & Compilazione per OS specifico. \\
Sicurezza & Sandboxed. & \texttt{unsafe}: un errore Rust può crashare il processo .NET. \\
Complessità & Richiede il pacchetto NuGet \texttt{Wasmtime}. & Nessuna dipendenza esterna. \\
\hline
\end{tabular}
\caption{Confronto tra le strategie di bridge WASM e Native FFI.}
\label{tab:wasm_vs_ffi}
\end{table}

La scelta è ricaduta sull'\textbf{approccio FFI nativo} per la priorità assegnata alle prestazioni massimali. In un contesto dove la verifica delle prove deve avvenire entro la finestra temporale di 12 secondi di uno slot Ethereum, ogni overhead è significativo. L'approccio WASM, pur offrendo un migliore isolamento, introduce un livello di indirezione (la VM WebAssembly) che penalizza il throughput in scenari di verifica intensiva. In cambio della perdita del sandboxing, il bridge nativo offre zero overhead di marshalling ed esecuzione diretta delle istruzioni sulla CPU del nodo.

\subsection{Il Bridge FFI: \texttt{native-zk-verifier}}\label{sec:ffi_bridge}

Il ponte tra i due runtime è realizzato dal crate Rust \texttt{native-zk-verifier}, configurato come \texttt{cdylib} (libreria condivisa C-compatibile) e compilato con Rust nightly, come richiesto dalle dipendenze crittografiche.

\subsubsection{Lato Rust}

Il crate espone una singola funzione C-ABI che funge da dispatcher verso il verifier appropriato:

\begin{minipage}{\textwidth}
\begin{lstlisting}[language=Rust, caption={Dispatcher FFI in Rust}, style=listingsStyle]
#[no_mangle]
pub extern "C" fn verify(
    zk_type: u32,
    proof_ptr: *const u8, proof_len: usize,
    vk_ptr: *const u8,    vk_len: usize
) -> i32
\end{lstlisting}
\end{minipage}
\\

Il parametro \texttt{zk\_type} seleziona il verifier tramite un enum intero. L'intera chiamata è avvolta in \texttt{panic::catch\_unwind}, un costrutto Rust che cattura eventuali panic nel codice dei verifier impedendo che si propaghino al processo .NET host, il quale crasherebbe irreversibilmente. I valori di ritorno seguono una convenzione semplice: \texttt{1} (Valid), \texttt{0} (Invalid), \texttt{-1} (Error).

Internamente, ciascun modulo verifier (\texttt{zisk.rs}, \texttt{openvm.rs}, \texttt{pico.rs}, \texttt{sp1\_hypercube.rs}, \texttt{airbender.rs}) implementa un trait comune \texttt{Verifier} con una funzione \texttt{verify(proof, vk) -> Result<bool>}, astraendo le differenze di serializzazione e di librerie sottostanti.

\subsubsection{Lato C\#}

Sul versante .NET, la dichiarazione P/Invoke utilizza il source-generated marshalling di .NET 8+:

\begin{minipage}{\textwidth}
\begin{lstlisting}[language={[Sharp]C}, caption={Dichiarazione P/Invoke in C\#}, style=listingsStyle]
[LibraryImport("native_zk_verifier")]
public static partial int verify(
    int zk_type,
    byte[] proof, nuint proof_len,
    nint vk_ptr,  nuint vk_len);
\end{lstlisting}
\end{minipage}
\\

La scelta di \texttt{LibraryImport} rispetto al più tradizionale \texttt{DllImport} elimina l'overhead del marshalling a runtime, generando il codice di interoperabilità staticamente in fase di compilazione.

\subsubsection{Gestione della Memoria}

Il modello di ownership della memoria tra i due runtime è un aspetto critico del bridge:

\begin{itemize}
    \item \textbf{Verification Key (VK):} Allocata in \textbf{memoria unmanaged} tramite \texttt{Marshal.AllocHGlobal}, al di fuori del Garbage Collector .NET. Questo evita che il GC sposti o liberi la VK durante la chiamata FFI. Il wrapper \texttt{ZkProofVerifier} implementa \texttt{IDisposable} per il rilascio deterministico della memoria unmanaged.
    \item \textbf{Proof data:} Owned dal GC di C\# e passata come array \texttt{byte[]}. Il runtime .NET effettua il pinning automatico durante la chiamata P/Invoke.
    \item \textbf{Rust:} Agisce come un ``guest'' puramente funzionale: riceve puntatori read-only, esegue la verifica, e ritorna un intero. Nessuna allocazione o liberazione di memoria avviene sul lato Rust.
\end{itemize}

\subsubsection{Build Automatica}

L'integrazione tra i due sistemi di build è gestita da un target MSBuild custom (\texttt{BuildRustDependency}) nel file \texttt{.csproj}:

\begin{enumerate}
    \item Esegue \texttt{cargo build --release} nella directory \texttt{native-zk-verifier/}.
    \item Copia la libreria nativa risultante (\texttt{.so}/\texttt{.dll}/\texttt{.dylib}) nella directory di output del progetto .NET.
\end{enumerate}

Questo garantisce che un singolo comando \texttt{dotnet build} compili automaticamente sia il codice C\# che il codice Rust, semplificando notevolmente il workflow di sviluppo.

\subsection{L'Applicazione Console Multi-Verifier}

Con il bridge FFI operativo, la prima fase culmina nella creazione di \texttt{EthProofValidator}: un'applicazione console standalone C\# che valida blocchi mainnet reali senza alcun overhead del client Nethermind.

\subsubsection{Architettura}

L'applicazione è strutturata in quattro componenti:

\begin{itemize}
    \item \texttt{Program.cs}: Entry point CLI che accetta \texttt{<LatestBlockId> <BlockCount>} e itera sui blocchi invocando il validatore.
    \item \texttt{BlockValidator}: Orchestratore che per ciascun blocco recupera le prove da ethproofs.org, le distribuisce ai verifier in parallelo (\texttt{Task.WhenAll}), e applica la logica di majority voting.
    \item \texttt{VerifierRegistry}: Registry pattern implementato con \\ \texttt{ConcurrentDictionary<clusterId, ZkProofVerifier>}. I verifier vengono creati \textit{lazily} al primo incontro con un nuovo cluster, scaricando la relativa VK dall'API di ethproofs.org.
    \item \texttt{EthProofsApiClient}: Client HTTP per le API REST di ethproofs.org (fetch prove, download binari, download VK).
\end{itemize}

\subsubsection{Majority Voting}

La logica di accettazione di un blocco segue un algoritmo di \textit{majority voting} multi-prover: un blocco è considerato valido se almeno il 50\% delle prove non-skipped risulta valido (\texttt{validCount * 2 >= totalCount}). Questa soglia rappresenta un compromesso tra sicurezza e tolleranza ai guasti: un singolo verifier con errori non invalida l'intero blocco, purché la maggioranza converga sullo stesso risultato.

\subsubsection{Risultati della Demo Standalone}

L'applicazione è stata eseguita su una sequenza di 50 blocchi consecutivi della mainnet Ethereum (\#24199951--\#24200000), con 5 prove da 5 zkVM distinte per ciascun blocco. I risultati chiave:

\begin{itemize}
    \item \textbf{Tutti i 50 blocchi validati con successo} (5/5 prove valide per blocco).
    \item \textbf{Tempi di verifica per singola prova:} $\sim$5--50\,ms (Zisk, Airbender), $\sim$90--140\,ms (OpenVM, Pico).
    \item \textbf{Tempo totale per blocco:} $\sim$1.6--2.5\,s, dominato dal download HTTP delle prove e non dalla verifica crittografica.
\end{itemize}

Di seguito è riportato un estratto dei log di esecuzione che illustra l'elaborazione di una sottosequenza di 3 blocchi:


\begin{minipage}{\textwidth}
\begin{lstlisting}[
  basicstyle=\ttfamily\small,
  backgroundcolor=\color{white},
  frame=single,
  rulecolor=\color{black!30},
  framesep=8pt,
  xleftmargin=5pt,
  xrightmargin=5pt,
  numbers=none,
  keepspaces=true,
  caption=Estratto log di esecuzione EthProofValidator,
  keywords={Valid, ACCEPTED},
  keywordstyle=\color{green!50!black}\bfseries,
  ndkeywords={Processing, Block, Total, Time},
  ndkeywordstyle=\color{blue!60!black}\bfseries
]
Processing Block #24199996
   Proof 5218300 - Zisk            : Valid (43 ms)
   Proof 5218298 - Zisk            : Valid (17 ms)
   Proof 5218303 - OpenVM          : Valid (114 ms)
   Proof 5218306 - Airbender       : Valid (13 ms)
   Proof 5218301 - Pico            : Valid (89 ms)
   --------------------------------
BLOCK #24199996 ACCEPTED (5/5)
Total Time: 2059 ms

Processing Block #24199997
   Proof 5218316 - Zisk            : Valid (40 ms)
   Proof 5218315 - Zisk            : Valid (15 ms)
   Proof 5218322 - OpenVM          : Valid (107 ms)
   Proof 5218326 - Airbender       : Valid (12 ms)
   Proof 5218319 - Pico            : Valid (112 ms)
   --------------------------------
BLOCK #24199997 ACCEPTED (5/5)
Total Time: 1868 ms

Processing Block #24199998
   Proof 5218337 - Zisk            : Valid (35 ms)
   Proof 5218340 - OpenVM          : Valid (111 ms)
   Proof 5218335 - Zisk            : Valid (6 ms)
   Proof 5218346 - Airbender       : Valid (33 ms)
   Proof 5218339 - Pico            : Valid (115 ms)
   --------------------------------
BLOCK #24199998 ACCEPTED (5/5)
Total Time: 1857 ms
\end{lstlisting}
\end{minipage}

Questi risultati dimostrano la fattibilità del modello multi-verifier e confermano che la verifica crittografica è trascurabile rispetto alla latenza di rete, validando empiricamente la scelta del bridge FFI nativo.

% ------------------------------------------------------------------
\section{Fase 2: Nethermind Integration}\label{sec:integration}

Con il verificatore standalone operativo, la seconda fase affronta l'integrazione della logica di validazione ZK nel ciclo di vita del client Nethermind, senza eseguire un fork del repository upstream.

\subsection{Alternative Considerate}

Prima di giungere alla soluzione finale, sono state valutate e scartate diverse strategie di integrazione:

\begin{enumerate}
    \item \textbf{Strategy Pattern nel \texttt{NewPayloadHandler}:} Definire un'interfaccia \texttt{IPayloadExecutionStrategy} con due implementazioni \\ (\texttt{LocalExecutionStrategy} per l'EVM standard, \\\texttt{ZkValidationStrategy} per la verifica ZK). Scartata perché richiedeva la modifica del codice core del gestore standard, violando il principio di non-invasività.

    \item \textbf{Modifica del \texttt{BlockchainProcessor}:} Come suggerito dal team Nethermind, intervenire a livello del processore di blocchi per bypassare l'esecuzione EVM. Scartata perché il \texttt{BlockchainProcessor} è fortemente accoppiato alla macchina a stati standard e avrebbe richiesto modifiche troppo invasive.

    \item \textbf{Override diretto del \texttt{MergePlugin}:} Estendere il plugin esistente che gestisce il consenso post-Merge. Scartata perché avrebbe creato una dipendenza fragile dalla struttura interna del \texttt{MergePlugin}, soggetta a rotture ad ogni aggiornamento del client.
\end{enumerate}

\subsection{L'Architettura a Plugin}

La soluzione adottata sfrutta il sistema di plugin modulare di Nethermind, basato sul container di Dependency Injection \textbf{Autofac}. L'interfaccia chiave è \texttt{IConsensusWrapperPlugin}, che permette di:

\begin{itemize}
    \item \textbf{Registrare tipi custom} nel container DI tramite un \texttt{Autofac.Module}.
    \item \textbf{Sovrascrivere i gestori} dell'Engine API senza modificare il codice core di Nethermind.
\end{itemize}

Questo approccio è il medesimo utilizzato dal \texttt{MergePlugin} (che gestisce il consenso post-Merge) e dallo \texttt{ShutterPlugin} (che estende la logica di consenso per il protocollo Shutter), garantendo compatibilità architetturale con il design del client\cite{nethermind_plugin}.

\subsection{\texttt{ZkValidationPlugin}: Il Punto di Ingresso}

Il plugin \texttt{ZkValidationPlugin} implementa \texttt{IConsensusWrapperPlugin} ed è il punto di ingresso dell'intero PoC. La sua attivazione è condizionale:

\begin{minipage}{\textwidth}
\begin{lstlisting}[language={[Sharp]C}, caption={Attivazione condizionale del plugin ZK}, style=listingsStyle]
public bool Enabled => zkConfig.Enabled && mergeConfig.Enabled;
\end{lstlisting}
\end{minipage}
\\

Il plugin si attiva solo quando entrambe le condizioni sono soddisfatte: la flag \texttt{ZkValidation.Enabled} è impostata a \texttt{true} nel file di configurazione, e il modulo \texttt{Merge} è attivo (requisito per il funzionamento dell'Engine API post-Merge).

Il \texttt{ZkValidationPluginModule} (Autofac Module) registra nel container DI tutti i servizi necessari: 
\begin{itemize}
    \item \texttt{ZkValidationService} (orchestratore condiviso)
    \item \texttt{ZkNewPayloadHandler}
    \item \texttt{ZkForkchoiceUpdatedHandler}
    \item \texttt{BlockValidator}
    \item \texttt{EthProofsApiClient}
\end{itemize}

\subsection{Override dell'Engine API}

\subsubsection{\texttt{ZkNewPayloadHandler}}

Il gestore ZK di \texttt{engine\_newPayload} sostituisce il \texttt{NewPayloadHandler} standard e semplifica radicalmente il pre-processing. La tabella \ref{tab:preprocessing} sintetizza i controlli mantenuti e quelli eliminati.
\\
\begin{table}[H]
\centering
\begin{tabular}{l c p{6cm}}
\hline
\textbf{Controllo} & \textbf{ZK} & \textbf{Motivazione} \\
\hline
Decodifica Blocco         & $\checkmark$ & Necessario per ottenere la struttura dati \\
Validazione Hash          & $\checkmark$ & Integrità crittografica dell'header \\
Invalid Chain Tracker     & $\checkmark$ & Propagazione di invalidità \\
\hline
Controllo Pre-Pivot       & $\times$ & Irrilevante senza sincronizzazione \\
Ricerca Header Genitore   & $\times$ & Nessun database locale \\
Inserimento Blocchi Orfani & $\times$ & Nessun sync in corso \\
Controllo ShouldProcessBlock & $\times$ & Si vuole sempre la validazione ZK \\
Controlli TTD             & $\times$ & Solo blocchi post-merge \\
\textbf{Esecuzione EVM}   & $\times$ & \textbf{Sostituito dalla verifica ZK} \\
\hline
\end{tabular}
\caption{Confronto dei controlli di pre-processing: gestore standard vs. gestore ZK.}
\label{tab:preprocessing}
\end{table}

Dopo il pre-processing semplificato, il gestore delega la validazione al \texttt{ZkValidationService}, il quale verifica in un tempo massimo la disponibilità delle prove, le scarica ed esegue la validazione locamente restituendo \texttt{VALID} o \texttt{INVALID} di conseguenza. Se le prove non sono disponibili l'EL ritornerà lo stato di \texttt{SYNCING}, aspettando nuovamente la chiamata di \texttt{engine\_newPayload} dal CL.

\subsubsection{\texttt{ZkForkchoiceUpdatedHandler}}

Il gestore ZK di \texttt{engine\_forkchoiceUpdated} bypassa il controllo \\\texttt{WasProcessed}, una flag normalmente impostata a \texttt{true} solo dopo l'esecuzione EVM e quindi sempre \texttt{false} in modalità ZK. Le operazioni residue sono:

\begin{enumerate}
    \item Verifica che il blocco non appartenga a una catena invalidata (\texttt{IsOnInvalidChain}).
    \item Ricerca del blocco nel \texttt{blockTree}.
    \item Aggiornamento del \texttt{blockTree.ForkChoiceUpdated(...)}.
    \item Restituzione di \texttt{VALID} con il \texttt{headBlockHash} aggiornato.
\end{enumerate}

% ------------------------------------------------------------------
\section{Fase 3: The ZK Validation Service}\label{sec:validation_service}

La terza fase integra la logica di verifica nel contesto operativo completo del client, affrontando le sfide di latenza dei prover, gestione della cache, e interazione con il protocollo di consenso.

\subsection{\texttt{ZkValidationService}: L'Orchestratore}

Il \texttt{ZkValidationService} è il servizio condiviso tra i due gestori dell'Engine API. Le sue responsabilità includono:

\begin{itemize}
    \item \textbf{Cache dei blocchi validati:} Prima di ogni verifica, il servizio controlla se il blocco è già presente nel \texttt{blockTree} tramite \texttt{FindBlock(hash)}. In caso affermativo, ritorna immediatamente \texttt{VALID} senza re-verifica, evitando lavoro ridondante quando il CL invia lo stesso blocco sia via \texttt{newPayload} che via \\\texttt{forkchoiceUpdated}.

    \item \textbf{Inserimento nel BlockTree:} Su risultato \texttt{VALID}, il blocco viene inserito con \\\texttt{blockTree.Insert(block, SaveHeader, BeaconBlockInsert)}, persistendo solo l'header nel \texttt{MemDb} senza scrivere alcun dato di stato.

    \item \textbf{Gestione dell'invalidità:} Su risultato \texttt{INVALID}, il servizio invoca \texttt{invalidChainTracker.OnInvalidBlock(hash, parentHash)}, marcando l'intera catena discendente come invalida e propagando l'informazione ai successivi blocchi figli.
\end{itemize}

\subsection{La Gestione del Protocollo SYNCING}\label{sec:syncing}

La sfida ingegneristica principale della Demo 1 (sezione \ref{sec:demo1}) è stata la gestione della latenza dei prover. Quando un blocco appena prodotto dalla rete raggiunge il nodo, la prova crittografica potrebbe non essere ancora disponibile su ethproofs.org (il tempo medio di proving è di circa 5--15 secondi). Senza un meccanismo di gestione adeguato, il nodo entrerebbe in uno stato \texttt{SYNCING} permanente, poiché il CL continuerebbe a inviare nuovi blocchi prima che i precedenti siano verificati.

La soluzione implementata è un \textbf{pattern Active-Wait}: il \\\texttt{ZkValidationService} esegue un retry loop configurabile che blocca il thread RPC durante l'attesa della prova. Il flusso operativo è il seguente:

\begin{enumerate}
    \item Il \texttt{BlockValidator} interroga ethproofs.org e non trova prove disponibili $\rightarrow$ ritorna \texttt{SYNCING}.
    \item Il servizio attende \texttt{RetryDelayMs} e ritenta la query fino al numero massimo di tentativi.
    \item Il CL (Lighthouse) riceve \texttt{SYNCING} e passa in modalità ``optimistic'' (\texttt{WARN Head is optimistic}).
    \item Una volta che la prova è disponibile e verificata, il servizio ritorna \texttt{VALID}.
    \item Lighthouse transita a \texttt{INFO Synced ... (verified)}.
\end{enumerate}

Questo approccio rappresenta un \textbf{workaround funzionale} per il PoC: il blocco del thread RPC per un massimo di $\sim$10 secondi è accettabile in un contesto di dimostrazione, ma non sarebbe sostenibile in un ambiente di produzione dove la reattività dell'Engine API è critica. Un'implementazione production-ready richiede prima un addattamento e modifica del protocollo attuale.

\subsection{Adattamento del BlockValidator}

Il \texttt{BlockValidator} sviluppato nella sezione \ref{sec:foundation} (standalone) è stato adattato al contesto del plugin Nethermind con le seguenti modifiche:

\begin{itemize}
    \item Integrazione con il sistema di logging di Nethermind (\texttt{ILogManager}) per una diagnostica unificata.
    \item Condivisione del \texttt{VerifierRegistry} e dell'\texttt{EthProofsApiClient} tramite il container DI di Autofac, anziché istanziazione diretta.
    \item Medesima logica di majority voting, con gestione degli errori adattata: un errore su un singolo verifier non interrompe il processo, ma viene registrato e conteggiato nella votazione.
\end{itemize}

% ------------------------------------------------------------------
\section{Profilo di Configurazione}\label{sec:config}

Il file \texttt{zk-mainnet.json} definisce il profilo operativo che trasforma Nethermind in un validatore ``Fort Mode''. La \ref{tab:config_disabled} riassume le componenti disabilitate e le relative motivazioni.
\\
\begin{table}[H]
\centering
\begin{tabular}{l l l p{5cm}}
\hline
\textbf{Namespace} & \textbf{Parametro} & \textbf{Valore} & \textbf{Motivazione} \\
\hline
Init & DiagnosticMode & MemDb & Tutte le store in memoria, zero I/O su disco \\
Init & DiscoveryEnabled & false & Nessun peer discovery P2P \\
Init & PeerManagerEnabled & false & Nessuna gestione connessioni P2P \\
Init & ProcessingEnabled & false & Disabilita il \texttt{BlockchainProcessor} \\
Sync & SynchronizationEnabled & false & Nessuna sincronizzazione \\
TxPool & Size & 0 & Nessun gossip di transazioni, nessun mempool \\
\hline
\end{tabular}
\caption{Moduli disabilitati nel profilo \texttt{zk-mainnet}.}
\label{tab:config_disabled}
\end{table}

I moduli che restano attivi sono esclusivamente quelli necessari al funzionamento del validatore stateless: \texttt{JsonRpc} (con i soli moduli \texttt{Net}, \texttt{Web3}, \texttt{Engine} per la comunicazione con il CL), \texttt{Merge} (prerequisito per l'Engine API post-Merge), \texttt{ZkValidation} (attivazione del plugin), \texttt{HealthChecks} e \texttt{Metrics} (monitoraggio e telemetria).

Il risultato è un nodo che non ha database su disco (MemDb al posto dei $\sim$2\,TB di RocksDB), non scopre né gestisce peer (riceve blocchi esclusivamente dal CL adiacente), non esegue transazioni (il \\ \texttt{BlockchainProcessor} è disattivato), non partecipa al gossip delle transazioni, e valida esclusivamente tramite prove ZK. Questo profilo materializza concretamente il concetto di ``Fort Mode'' descritto nel Capitolo~\ref{chap:introduzione}: un validatore leggero, deployabile su hardware consumer, che contribuisce alla sicurezza della rete delegando la generazione delle prove ai ``Beast'' e concentrandosi sulla sola verifica crittografica.



% ============================================================
%  CAPITOLO 5: VALIDAZIONE
% ============================================================
\chapter{Validazione}\label{chap:validazione}

La validazione del PoC ZeroNethermind si articola attorno a due domande di ricerca complementari:

\begin{enumerate}
    \item \textbf{Fattibilità operativa:} Un verificatore stateless può mantenere la sincronizzazione con la chain L1 Mainnet in tempo reale, operando esclusivamente tramite prove crittografiche?
    \item \textbf{Impatto sulle risorse:} Quali sono le differenze quantitative in termini di velocità di verifica, consumo di CPU, RAM, I/O su disco e storage tra un nodo standard a ri-esecuzione completa e un nodo ZeroNethermind?
\end{enumerate}

Per rispondere a queste domande, sono state progettate due campagne sperimentali distinte: la \textbf{Demo~1} (Sezione \ref{sec:demo1}), che effettua una validazione live sulla Mainnet Ethereum tramite prove fornite da \texttt{ethproofs.org}, e la \textbf{Demo~2} (Sezione \ref{sec:demo2}, che utilizza un ambiente controllato basato su Kurtosis\cite{kurtosis} per isolare e confrontare le metriche hardware in condizioni deterministiche.


\section{Metodologia}\label{sec:metodologia}

\subsection{Ambiente di Test}\label{sec:ambiente_test}

Tutti gli esperimenti sono stati eseguiti sulla medesima macchina fisica, utilizzando container Docker per garantire l'isolamento e la riproducibilità. Le specifiche hardware sono riportate nella tabella \ref{tab:hardware}.
\\
\begin{table}[H]
\centering
\begin{tabular}{l l}
\hline
\textbf{Componente} & \textbf{Specifica} \\
\hline
Sistema Operativo & Manjaro Linux, kernel 6.12.73 (64-bit) \\
CPU & 4 x Intel Core i5-7400 @ 3.00\,GHz \\
RAM & 16\,GiB DDR4 \\
GPU & NVIDIA GeForce GTX 1060 6\,GB \\
Rete & Connessione domestica (FTTC) \\
Containerizzazione & Docker Engine con \texttt{docker-compose} \\
\hline
\end{tabular}
\caption{Specifiche hardware della macchina di test.}
\label{tab:hardware}
\end{table}

La scelta di hardware consumer di fascia media è intenzionale: l'obiettivo del PoC è dimostrare che la validazione ZK è realizzabile su dispositivi comuni, senza necessità di server specializzati.


\subsection{Demo~1: Validazione in Tempo Reale sulla Mainnet}\label{sec:demo1}

\subsubsection{Setup}

La Demo~1 consiste nel deployment di ZeroNethermind accoppiato con un Consensus Layer Lighthouse, entrambi eseguiti come container Docker sulla macchina locale.

\begin{itemize}
    \item \textbf{Execution Layer:} ZeroNethermind con il profilo \texttt{zk-mainnet} (sezione \ref{sec:config}), configurato in modalità \texttt{MemDb} (zero storage su disco), con discovery P2P, peer manager, sincronizzazione e \\\texttt{BlockchainProcessor} disabilitati.
    \item \textbf{Consensus Layer:} Lighthouse in modalità \texttt{checkpoint-sync}, sincronizzato al tip della Beacon Chain tramite il checkpoint pubblico \texttt{mainnet.checkpoint.sigp.io}.
    \item \textbf{Fonte delle prove:} \texttt{ethproofs.org} via HTTP API, interrogato automaticamente dal \texttt{BlockValidator} per ogni blocco ricevuto tramite \texttt{engine\_newPayload}.
    \item \textbf{Monitoraggio:} Prometheus\cite{prometheus} + cAdvisor\cite{cAdvisor} per il campionamento delle metriche container (CPU, RAM, I/O, bandwidth) con frequenza di 15 secondi.
\end{itemize}

\subsubsection{Obiettivo}

Dimostrare la fattibilità operativa del paradigma ``Fort Mode'': il nodo, privo di database e di capacità di esecuzione EVM, deve riuscire a validare i blocchi della Mainnet in tempo reale, mantenendo la sincronizzazione con il CL e rispondendo a ciascuna chiamata \texttt{engine\_newPayload} con uno status \texttt{VALID} entro i vincoli temporali dello slot.

\subsubsection{Procedura}

Il test è stato condotto su una finestra di circa 17 minuti, durante i quali il nodo ha ricevuto e validato \textbf{121~blocchi Mainnet} (nel range \texttt{\#24628200} -- \texttt{\#24628397}), corrispondenti a circa 3 epoche della Beacon Chain. Per ciascun blocco, il \texttt{BlockValidator} ha:

\begin{enumerate}
    \item Interrogato \texttt{ethproofs.org} per ottenere l'elenco delle prove disponibili.
    \item Scaricato in parallelo le prove binarie generate da diversi zkVM indipendenti: \textbf{Zisk} (tipicamente 2 istanze), \textbf{OpenVM}, \textbf{Pico} e \textbf{Airbender}.
    \item Invocato il bridge FFI nativo (\texttt{libnative\_zk\_verifier.so}) per la verifica crittografica di ciascuna prova.
    \item Eseguito il majority voting: il blocco è accettato se almeno il 50\% delle prove verificate risulta valido.
    \item Inserito il blocco (\textit{header-only}) nel \texttt{BlockTree} in-memory e restituito \texttt{VALID} al CL.
\end{enumerate}

Per ogni blocco validato, il sistema ha registrato nel log: il tempo totale di validazione, il tempo di download massimo (bottleneck del parallelismo HTTP), il tempo di verifica FFI massimo (bottleneck del parallelismo crittografico), e il conteggio prove valide/totali. I dati sono stati estratti dai log del container e aggregati in formato CSV.


\subsection{Demo~2: Ambiente Controllato (Kurtosis Lab)}\label{sec:demo2}

\subsubsection{Setup}

La Demo~2 utilizza un'enclave Kurtosis per creare una devnet Ethereum locale completamente isolata, configurata per confrontare direttamente il footprint hardware di un nodo Nethermind standard e di un nodo ZeroNethermind.

\begin{itemize}
    \item \textbf{Nodi standard:} 2 istanze Nethermind con configurazione di default, equipaggiate con validatori e MEV-Boost.
    \item \textbf{Nodo ZeroNethermind:} 1 istanza con il profilo \texttt{zk-mainnet} e il plugin \texttt{ZkValidation} attivo, connessa al \textbf{Mock Prover} locale.
    \item \textbf{Consensus Layer:} Lighthouse per tutti i nodi.
    \item \textbf{Mock Prover:} Un servizio containerizzato (\texttt{zero-prover}) che replica l'API di \texttt{ethproofs.org}. Per ogni blocco restituisce \textbf{5 istanze} della medesima prova crittografica Airbender pre-generata e valida, simulando il carico di lavoro reale in cui il nodo scarica e verifica prove da 5 zkVM indipendenti in parallelo. Questa scelta consente di misurare il tempo di verifica FFI in condizioni realistiche, eliminando al contempo la variabilità della latenza di rete verso un servizio esterno.
    \item \textbf{Carico di lavoro:} \texttt{spamoor} genera un flusso continuo di transazioni per saturare i blocchi al Gas Limit corrente.
    \item \textbf{Monitoraggio:} Prometheus, Grafana e \texttt{ethereum\_metrics\_exporter} per la raccolta integrata delle metriche di tutti i nodi.
\end{itemize}

L'utilizzo di MEV-Boost per entrambe le tipologie di nodo è intenzionale: delegando la costruzione dei blocchi a un builder esterno, tutti i nodi si limitano alla sola fase di validazione, che è esattamente il componente sotto test.

\subsubsection{Obiettivo}

Quantificare il divario di risorse hardware tra validazione per ri-esecuzione EVM ($O(N)$) e validazione per verifica crittografica ($O(1)$), in un ambiente dove le condizioni di rete, il carico transazionale e l'hardware sono identici per tutti i nodi.

\subsubsection{Scenari di Test}

La Demo~2 prevede tre scenari distinti, ciascuno progettato per isolare una specifica dimensione di confronto:

\begin{enumerate}
    \item \textbf{Scenario~1 --- Resource Footprint:} Entrambi i tipi di nodo validano la stessa catena con blocchi saturi (30M Gas) generati da \texttt{spamoor}. Si confrontano CPU, RAM, I/O e storage su una finestra di $N$ epoche.
    \item \textbf{Scenario~2 --- Sync Time:} Si avvia la devnet con soli nodi standard per $N$ minuti (accumulando storia), poi si iniettano simultaneamente un nodo Nethermind fresco e un nodo ZeroNethermind, misurando il tempo necessario a ciascuno per raggiungere il tip.
    \item \textbf{Scenario~3 --- Scalabilità del Gas Limit:} Si eseguono sessioni separate con valori crescenti di \texttt{genesis\_gas\_limit} (30M, 60M, 120M, 250M) per tracciare la curva $O(N)$ della ri-esecuzione EVM contro la costante $O(1)$ della verifica ZK.
\end{enumerate}

I risultati della Demo~2 sono presentati nella sezione \ref{sec:demo2_risultati}.


\section{Demo~1: Validazione sulla Mainnet}\label{sec:demo1_risultati}

I risultati della Demo~1 sono presentati attraverso due visualizzazioni complementari:

\begin{minipage}{\textwidth}
\begin{figure}[H]
    \makebox[\textwidth][c]{\includegraphics[width=1.1\linewidth]{img/graph1.png}}
    \caption{Composizione del Tempo Totale di Validazione}
    \label{fig:demo1_graph1}
\end{figure}
\end{minipage}\\

Il \textbf{Grafico \ref{fig:demo1_graph1}} illustra la scomposizione del tempo totale per validare ciascun blocco all'interno del gestore \texttt{engine\_newPayload}. Come si evince chiaramente dalla predominanza della componente blu rispetto a quella verde, il \textit{Max Download Time} (il tempo impiegato per recuperare le prove tramite HTTP da ethproofs.org) costituisce il vero e unico collo di bottiglia del processo, assorbendo regolarmente oltre il 90\% del tempo totale. Al contrario, il \textit{Max Verify Time} (in verde), ovvero il tempo di calcolo effettivo speso dalla CPU per la verifica crittografica tramite il bridge FFI, si mantiene costante e irrisorio (nell'ordine dei 100-150 ms), indipendentemente dalla complessità del blocco. I picchi anomali visibili (fino a $\sim$9000 ms) sono imputabili esclusivamente a latenze di rete transitorie del provider esterno.

\begin{minipage}{\textwidth}
\begin{figure}[H]
    \makebox[\textwidth][c]{\includegraphics[width=1.1\linewidth]{img/graph2.png}}
    \caption{Analisi Prestazionale: Nodo Stateless ZeroNethermind}
    \label{fig:demo1_graph2}
\end{figure}
\end{minipage}\\

Il \textbf{Grafico \ref{fig:demo1_graph2}} conferma empiricamente il successo del paradigma "Zero State" attraverso l'analisi temporale dell'utilizzo delle risorse hardware. Nel pannello superiore, l'utilizzo della memoria RAM (linea blu) si mantiene straordinariamente basso e piatto, assestandosi costantemente al di sotto dei 400 MB. Questo dimostra l'efficacia dell'eliminazione del database di stato persistente e della Mempool. Nel pannello inferiore, la metrica più significativa è il traffico I/O su disco (linea viola), che, dopo un minimo assorbimento in fase di avvio, si azzera del tutto ($0$ MB/s), confermando l'assenza totale delle pesanti operazioni di lettura/scrittura sull'albero MPT tipiche dell'EVM. Contestualmente, la banda di rete in ricezione (linea verde) mostra pattern regolari associati al download continuo delle prove crittografiche per ogni nuovo slot, confermando che il sistema si è spostato da un limite \textit{I/O-bound} a un limite \textit{Bandwidth-bound}.


La tabella \ref{tab:demo1_summary} riporta le statistiche aggregate della sessione.
\\
\begin{table}[H]
\centering
\begin{tabular}{l r r r r}
\hline
\textbf{Metrica} & \textbf{Media} & \textbf{Mediana} & \textbf{Min} & \textbf{Max} \\
\hline
Tempo totale di validazione (ms) & 2\,237 & 1\,950 & 1\,241 & 9\,400 \\
Max download time (ms)           & 1\,768 & 1\,447 & 817   & 8\,755 \\
Max verify time (ms)             & 122   & 116   & 48    & 301 \\
\hline
RAM container (MB)               & 360   & 335   & 272   & 457 \\
CPU container (\%, per-core)     & 57,7  & ---   & 32,2  & 74,8 \\
Bandwidth RX (KB/s)              & 333   & 295   & $<$1  & 1\,289 \\
I/O read (KB/s)                  & 77    & 34    & $<$1  & 476 \\
\hline
\end{tabular}
\caption{Statistiche aggregate della Demo~1 (121 blocchi Mainnet).}
\label{tab:demo1_summary}
\end{table}

Il risultato più significativo è la sproporzione tra il tempo di download e il tempo di verifica: il tempo medio di verifica crittografica FFI è di soli \textbf{122\,ms} (mediana 116\,ms), rappresentando appena il \textbf{5,5\%} del tempo totale medio di validazione. Il restante 94,5\% è dominato dalla latenza di rete per il download delle prove da \texttt{ethproofs.org}. Questo dato conferma che la verifica ZK è computazionalmente trascurabile rispetto al trasporto dei dati, e anticipa una proprietà fondamentale: quando le prove saranno incluse direttamente nel payload del blocco, il tempo di validazione si ridurrà esclusivamente alla componente di verifica, nell'ordine di $\sim$100\,ms.


\section{Demo~2: Analisi in Ambiente Controllato}\label{sec:demo2_risultati}

\subsection{Velocità di Verifica}\label{sec:velocita}

% TODO: Inserire Grafico --- Velocità di Verifica vs Esecuzione (Scenario~1).

% TODO: Inserire Grafico --- Impatto dell'innalzamento del Gas Limit (Scenario~3).
% Asse X: Gas Limit (30M -- 1 GigaGas), Asse Y: tempo di elaborazione (ms).
% La curva Nethermind cresce linearmente O(N), la curva ZeroNethermind resta piatta O(1).


\subsection{Footprint delle Risorse}\label{sec:footprint}

% TODO: Inserire Grafico Radar --- Confronto CPU, RAM, I/O, Bandwidth tra Nethermind e ZeroNethermind (Scenario~1).

% TODO: Inserire Grafico Linee --- Crescita dello Storage nel tempo (Scenario~1).
% Nethermind accumula stato su disco, ZeroNethermind resta a ~0 GB.


\subsection{Tempo di Sincronizzazione}\label{sec:sync_time}

% TODO: Inserire Grafico Barre --- Sync Time (Scenario~2).
% Confronto del tempo necessario per raggiungere il tip della catena.


\section{Discussione dei Risultati}\label{sec:discussione}

% TODO: Sintesi cross-demo (Demo 1 + Demo 2), osservazioni qualitative:
% stabilità del sync con Lighthouse, affidabilità del bridge FFI, multi-prover diversity.


Infine la tabella \ref{tab:confronto_finale} mostra le differenze architetturali e prestazionali tra un nodo Nethermind Full Node e il PoC ZeroNethermind.

\begin{table}[H]
\centering
\begin{tabular}{l p{5cm} p{5cm}}
\hline
\textbf{Dimensione} & \textbf{Nethermind Full Node} & \textbf{ZeroNethermind} \\
\hline
Trust Model & Ri-esecuzione locale di ogni transazione & Verifica crittografica \\
Storage & $\sim$1\,TB & $\sim$0\,GB (MemDb in RAM) \\
RAM & 16--32\,GB & $<$500\,MB \\
CPU & Saturazione multi-core durante l'esecuzione EVM & Picchi brevi ($\sim$120\,ms) solo durante la verifica FFI \\
I/O & Elevato (migliaia di IOPS su SSD) & Trascurabile ($<$1\,KB/s medio) \\
Bottleneck & I/O-bound (accesso a stato MPT) & Latency-bound (download proofs) \\
Validation & $O(N)$ --- lineare nel numero di transazioni & $O(1)$ --- costante \\
Sync Time & Ore/giorni & Secondi \\
RPC & Completa & Nessuna (attester-only) \\
Functions & Block building, serving, validazione & Solo validazione (Fort Mode) \\
\hline
\end{tabular}
\caption{Confronto tra Nethermind Full Node e ZeroNethermind Stateless Verifier.}
\label{tab:confronto_finale}
\end{table}


% ============================================================
%  CAPITOLO 6: CONCLUSIONI
% ============================================================
\chapter{Conclusioni}\label{chap:conclusioni}

\section{Riepilogo}
% Dimostrazione che un Execution Client può operare
% in modalità stateless affidandosi alla matematica anziché allo potenza di calcolo.

\section{Limitazioni}
% mock prover, latenza ethproofs, assenza di proving locale

\section{Lavori futuri}
 % EIP-8025 adoption, real-time proving integrato, proving on-chain, mobile deployment

\section{Considerazioni finali}
% impatto sulla roadmap Ethereum, passo verso "Lean ETH" e validatori mobile-native.


\renewcommand{\bibsection}{}
\chapter*{Riferimenti bibliografici}
\addcontentsline{toc}{chapter}{Riferimenti bibliografici}
\bibliography{refs}

\end{document}